% =====================================================================
%  main.tex — 主文件（Lead 独占）
%  编译：在 源码\ 目录下 `xelatex main.tex` 两遍（或 ..\工程文件\build.ps1）
%  注意：原文节号是 0,1,2,3,4，故正文前 \setcounter{section}{-1}
% =====================================================================
\documentclass[UTF8,a4paper,11pt]{ctexart}
% =====================================================================
%  src/preamble.tex  —  导言区（SY05-24 Gadde《超共形指标讲义》）
%  由 Lead 维护；译者只读。缺宏时在文末「追加区」append \newcommand，
%  并在回报里列出新增宏名（不得改动上面的已有定义）。
% =====================================================================
\usepackage[a4paper,margin=2.6cm,top=3.0cm,bottom=2.8cm]{geometry}
\usepackage{amsmath,amssymb,amsthm,mathtools}
\usepackage{mathrsfs}
\usepackage{graphicx}
\usepackage{array,booktabs,longtable,multirow}
\usepackage{caption}
\usepackage{enumitem}
\usepackage{xcolor}
\usepackage{fancyhdr}

% ---------- 字体（Windows，显式路径） ----------
\setCJKmainfont[Path=C:/Windows/Fonts/,BoldFont=simhei.ttf,ItalicFont=simkai.ttf]{simsun.ttc}
\setCJKsansfont[Path=C:/Windows/Fonts/]{simhei.ttf}
\setCJKmonofont[Path=C:/Windows/Fonts/]{simfang.ttf}
\setCJKfamilyfont{zhkai}[Path=C:/Windows/Fonts/]{simkai.ttf}
\setCJKfamilyfont{zhhei}[Path=C:/Windows/Fonts/]{simhei.ttf}
\setCJKfamilyfont{zhfang}[Path=C:/Windows/Fonts/]{simfang.ttf}
\providecommand{\kaishu}{\CJKfamily{zhkai}}
\providecommand{\heiti}{\CJKfamily{zhhei}}
\providecommand{\fangsong}{\CJKfamily{zhfang}}

% ---------- 数学宏（块文件里只许用这些，禁止自造） ----------
\newcommand{\sN}{\mathcal{N}}          % 超对称数 \sN=1, \sN=2, \sN=4
\newcommand{\sI}{\mathcal{I}}          % 指标（index）
\newcommand{\sO}{\mathcal{O}}          % 算符
\newcommand{\sH}{\mathcal{H}}          % 希尔伯特空间
\newcommand{\sZ}{\mathcal{Z}}          % 配分函数
\newcommand{\sL}{\mathcal{L}}          % 拉氏量
\newcommand{\sQ}{\mathcal{Q}}          % 超荷
\newcommand{\sU}{\mathcal{U}}
\newcommand{\sV}{\mathcal{V}}
\DeclareMathOperator{\Tr}{Tr}
\DeclareMathOperator{\PE}{PE}
\DeclareMathOperator{\Sym}{Sym}
\DeclareMathOperator{\adj}{adj}
\DeclareMathOperator{\fund}{fund}
\DeclareMathOperator{\sgn}{sgn}
\DeclareMathOperator{\Li}{Li}
\DeclareMathOperator{\SU}{SU}
\DeclareMathOperator{\SO}{SO}
\DeclareMathOperator{\Sp}{Sp}
\DeclareMathOperator{\USp}{USp}
\DeclareMathOperator{\U}{U}
\DeclareMathOperator{\GL}{GL}
\DeclareMathOperator{\diag}{diag}
\newcommand{\R}{\mathbb{R}}
\newcommand{\C}{\mathbb{C}}
\newcommand{\Z}{\mathbb{Z}}
\newcommand{\Q}{\mathbb{Q}}
\newcommand{\bbZ}{Z}
\newcommand{\eps}{\varepsilon}
\newcommand{\dd}{\mathrm{d}}
\newcommand{\ee}{\mathrm{e}}
\newcommand{\ii}{\mathrm{i}}
\newcommand{\ol}[1]{\overline{#1}}
\newcommand{\br}[1]{\bar{#1}}
\newcommand{\dg}{\dagger}
\newcommand{\half}{\tfrac{1}{2}}
\newcommand{\third}{\tfrac{1}{3}}
\newcommand{\twothird}{\tfrac{2}{3}}
% 表格里常用的列类型
\newcolumntype{C}[1]{>{\centering\arraybackslash}p{#1}}

% ---------- 定理类环境（编号一律手工给：\begin{exercise}[2.1]） ----------
\theoremstyle{definition}
\newtheorem*{exercise}{习题}
\newtheorem*{definition}{定义}
\newtheorem*{example}{例}
\newtheorem*{theorem}{定理}
\newtheorem*{proposition}{命题}
\newtheorem*{lemma}{引理}
\newtheorem*{corollary}{推论}
\newtheorem*{remark}{注}
\newtheorem*{solution}{解}

% 让 \begin{exercise}[2.1] 印成「习题 2.1」（无括号）
\makeatletter
\def\thmhead@plain#1#2#3{%
  \thmname{#1}\thmnumber{\@ifnotempty{#2}{ #2}}%
  \thmnote{\@ifnotempty{#3}{\hspace{0.35em}{#3}}}}
\let\thmhead\thmhead@plain
\makeatother

% ---------- 版式 ----------
\setlength{\parindent}{2em}
\linespread{1.18}
\setlist{nosep,leftmargin=2.2em}
\renewcommand{\arraystretch}{1.15}

\pagestyle{fancy}
\fancyhf{}
\fancyhead[L]{\small\kaishu 超共形指标讲义}
\fancyhead[R]{\small\thepage}
\renewcommand{\headrulewidth}{0.4pt}
\fancypagestyle{plain}{\fancyhf{}\fancyfoot[C]{\small\thepage}\renewcommand{\headrulewidth}{0pt}}

\ctexset{
  section = {
    format = \Large\bfseries\heiti,
    beforeskip = 2.2ex plus 1ex minus .2ex,
    afterskip  = 1.4ex plus .2ex,
    name = {},
    number = \arabic{section},
  },
  subsection = {
    format = \large\bfseries\heiti,
    beforeskip = 1.8ex plus .6ex minus .2ex,
    afterskip  = 1.0ex plus .2ex,
  },
}

\usepackage[hidelinks,bookmarksnumbered=false,bookmarksopen=true,
            pdftitle={超共形指标讲义（中文译注本）},
            pdfauthor={Abhijit Gadde}]{hyperref}
\pdfstringdefDisableCommands{\def\ol#1{#1}}

% ===================== 追加区（译者可 append \newcommand） =====================


\begin{document}

% ---------------------------- 封面 ----------------------------
\begin{titlepage}
\centering
\vspace*{3.0cm}
{\heiti\zihao{2} 超共形指标讲义\par}
\vspace{1.0cm}
{\kaishu\zihao{3} Lectures on the Superconformal Index\par}
\vspace{2.0cm}
{\zihao{4} Abhijit Gadde\par}
\vspace{0.5cm}
{\zihao{5} Department of Theoretical Physics,\\ Tata Institute of Fundamental Research, Mumbai 400005\par}
\vspace{2.0cm}
{\kaishu\zihao{4} 中文译注本\par}
\vspace{0.8cm}
{\zihao{5} 底本：arXiv:2006.13630v1 [hep-th]，2020 年 6 月 24 日，32 页\par}
\vfill
{\kaishu\zihao{5}
本译注本依据 arXiv:2006.13630v1（2020 年 YRISW 冬季学校讲义，拟刊于 J. Phys. A 专辑）全文翻译。\\
原文全部 110 条编号公式、8 个编号习题、2 张表格均与原文逐一对应；\\
原文笔误照译，并以注释指出正确形式。\par}
\vspace{1.0cm}
\end{titlepage}

% ---------------------------- 目录 ----------------------------
\tableofcontents
\newpage

% ------------------------- 译者说明 -------------------------
% translator_note.tex — 译者说明（放在目录之后、正文之前）
% 要求：4–6 段中文，说明 ①底本（YRISW 2020 冬季学校讲义，拟刊 J. Phys. A 专辑）与作者
% ②本译注本体例（编号公式手工 \tag、编号习题、两张表格复刻、参考文献不译、Tutor Note 不进排版）
% ③物理术语处理（指标/配分函数/轻子数…以 glossary.tsv 为准）④致谢与免责（仅供学习研究）。
\section*{译者说明}
\addcontentsline{toc}{section}{译者说明}


% ---------------------------- 摘要 ----------------------------
% abstract.tex — c01｜原文 p1｜论文标题页上的 Abstract（摘要）
% 0 条编号公式。
% [原文 p.1]
\noindent{\heiti\Large 摘要}\par
\addcontentsline{toc}{section}{摘要}
\vspace{0.4em}

在这些讲义中，我们给出超共形指标（superconformal index）的一个教学式导引。本稿是 2020 年冬季学校「YRISW 2020」上所作报告的文字版，将发表于 JPhysA 的一个专辑。讲义处于基础层次，面向有兴趣研究超共形指标的初年级研究生。

% [Tutor Note: 「教学式导引」对应 pedagogical introduction；作者有意面向初学者，正文保留大量口语与第一人称，译文不「润色」掉这种语气。]
% [Tutor Note: 术语 superconformal index 在 glossary 中冻结为「超共形指标」，不要写作「超共形指数」，以免与指数函数混淆。]
% [Tutor Note: YRISW 2020 是这次冬季学校的简称，原文未给全称；本译注本不擅自补全。]
% [Tutor Note: JPhysA 即 Journal of Physics A: Mathematical and Theoretical，原文说本稿将刊于其专辑。]
% [Tutor Note: 底本为 arXiv:2006.13630v1（2020 年 6 月 24 日，32 页）；本译本页标记用 PDF 页序，故正文第 1 页是 p.2。]
% [Tutor Note: 摘要只出现在 JHEP 标题页；正文 §0 的「讲义提纲」是另一段文字，二者不要混为一谈。]
% [Tutor Note: 原文 "interested in working with the superconformal index" 强调动手计算，译文用「研究」略失分寸，可理解为「以上手使用为目标」。]
% [Tutor Note: 摘要与 §0 都不带编号；§0 在原文的节计数器中算第 0 节，main.tex 用 \setcounter{section}{-1} 实现。]


% ---------------------------- 正文 ----------------------------
\setcounter{section}{-1}

\section{讲义提纲}
% sec0.tex — c01｜原文 p2｜§0 Plan of lectures（讲义提纲）
% 0 条编号公式；无编号环境。
% [原文 p.2]
超共形指标是超共形场论在径向量子化下的 Witten 指标。我们将首先在第 1 节结合超对称量子力学介绍 Witten 指标。除了理解指标受保护背后的基本原理，我们还将发展许多工具，特别是齐次指数化（plethystic exponentiation）；这些工具对于在更简单的设定下计算规范理论的超共形指标是必需的。在第 2 节，我们考虑超共形代数并引入超共形指标。基本组元——即手征多重态与矢量多重态——的指标将被计算出来。随后我们讨论指标在超共形代数短多重态重组下的不变性。我们的结论是：超共形指标在重组的意义下（up to recombinations）捕捉了理论的全部谱信息。在第 3 节，我们把超共形指标应用于所谓的「class S」对偶。在解释了具有超对称性增强的指标极限之后，我们描述与指标相伴的拓扑场论。最后，在第 4 节，我们计算大 $N$ 极限下 $\sN=4$ 的 $\SU(N)$ 超 Yang–Mills 的指标。先验地，$AdS_5$ 中的超对称黑洞被认为会对指标贡献 $O(N^2)$ 个态。
% [原文 p.3]
但我们发现，指标是 $O(1)$ 的，并且只由体（bulk）中的超引力态所饱和。我们对这一表面上的不匹配加以评论，并介绍最近一些从指标恢复黑洞熵的工作。指标的其它一些方面，例如它的模性质 [1, 2] 以及它与反常的联系 [3]，将不在这些讲义中讨论。我们只讨论 $S^3$ 上的指标。关于透镜空间（lens space）上指标的讨论，见 [4–7]。关于超共形指标的其它一些优秀综述，见 [8, 9]。

% [Tutor Note: 「径向量子化」把 $S^{d-1}\times\R$ 上的理论化为以径向为时间的哈密顿框架；超共形指标正是该框架下的 Witten 指标。]
% [Tutor Note: 「受保护」（protected）指该量在保持超对称的形变下不变，因而可在弱耦合自由极限精确计算。]
% [Tutor Note: 齐次指数把单字母谱提升为多粒子谱；它是第 1 节的技术核心，后文 (2.6)、(2.13) 会反复使用。]
% [Tutor Note: 「重组」（recombination）指短多重态合并成长多重态；指标对重组不变，故只捕捉「重组意义下」的谱信息。]
% [Tutor Note: 原文 class S 带引号且用花体 S；glossary 冻结为 class S 原样保留，可称「S 类理论」。]
% [Tutor Note: $AdS_5$ 黑洞的 $O(N^2)$ 微观态与指标 $O(1)$ 的落差是 §4 的主题；此处先给结论，细节见 (4.16) 附近。]
% [Tutor Note: 「透镜空间」（lens space）上的指标涉及更一般的边界条件；本讲义只讨论 $S^3$。]
% [Tutor Note: 本节是全篇路线图：可先跳到 §1 的谐振子例子建立直觉，再回看 §2 的公理化定义。]


\section{Witten 指标}
% sec1a.tex — c02｜原文 p3–5｜§1 Witten index + §1.1 Supersymmetric harmonic oscillator
% 编号公式 (1.1)–(1.17) 共 17 条；无编号环境。
% [原文 p.3]
\subsection{超对称谐振子}

考虑一个超对称谐振子
\[
L = \frac{1}{2}\dot x^2 - \frac{1}{2}x^2 + \ii\bar\psi\dot\psi - \bar\psi\psi .
\tag{1.1}
\]
该系统的哈密顿量为
\[
H = \dot\psi\,\Pi_\psi + \dot{\bar\psi}\,\Pi_{\bar\psi} + \dot x\,p - L
\tag{1.2}
\]
\[
\phantom{H} = \frac{1}{2}p^2 + \frac{1}{2}x^2 + \bar\psi\psi .
\tag{1.3}
\]
其中 $\Pi_\psi$、$\Pi_{\bar\psi}$ 与 $p$ 分别是 $\psi$、$\bar\psi$ 与 $x$ 的正则共轭量。对系统作量子化，得到产生与湮灭算符
\[
a_B^\dagger = \frac{1}{\sqrt2}(-\ii p + x),\qquad
a_B = \frac{1}{\sqrt2}(\ii p + x),\qquad
a_F^\dagger = \bar\psi,\qquad
a_F = \psi ,
\tag{1.4}
\]
它们满足 $[a_B, a_B^\dagger] = 1$ 与 $\{a_F, a_F^\dagger\} = 1$。用这些算符表出的哈密顿量为
\[
H = (a_B^\dagger a_B + a_F^\dagger a_F) .
\tag{1.5}
\]
不难看出，这个哈密顿量具有超对称性：$Q = a_B^\dagger a_F$，$\bar Q = a_F^\dagger a_B$。这些荷构成的代数为
\[
[Q, H] = [\bar Q, H] = 0 ,\qquad \{Q, \bar Q\} = H .
\tag{1.6}
\]
这是超对称代数的一个普遍特征。除这些连续对称性之外，理论还有一个与我们的讨论相关的离散 $\Z_2$ 对称性，即费米子数 $F$。算符 $F$ 数出态中费米子的数目（模 2）。

谱很容易算出。由于基态满足 $H|0\rangle = 0$，它同时也满足 $Q|0\rangle = 0$。激发态由产生算符作用在真空上得到。
\[
|\chi_B, n\rangle \equiv (a_B^\dagger)^n |0\rangle ,\qquad
|\chi_F, m\rangle \equiv (a_B^\dagger)^m a_F^\dagger |0\rangle .
\tag{1.7}
\]
% [原文 p.4]
第一组态的 $F = 0$，第二组态的 $F = 1$。配分函数也很容易算出：
\[
Z(x) = \Tr x^H = 1 + (x + x^2 + \cdots) + (x + x^2 + \cdots) = \frac{1+x}{1-x} ,\qquad x = \ee^{-\beta} .
\tag{1.8}
\]
让我们详细解释一个可用于计算自由理论配分函数的技巧。如上所述，谱完全由所有产生算符作用在真空上生成。这些产生算符彼此对易，因而独立地作用。方便的作法是在「单个字母」上计算配分函数，也就是在只含单个粒子的态上计算。这称为单字母配分函数，记作 $z(x)$。分别对玻色子与费米子计算它是有用的，二者分别记作 $z_B$ 与 $z_F$。在此情形下，只有一个玻色产生算符和一个费米产生算符。
\[
\begin{aligned}
z(x) &= z_B(x) + z_F(x),\\
z_B(x) &= \Tr_{\text{bosonic letters}} x^H = x,\\
z_F(x) &= \Tr_{\text{fermionic letters}} x^H = x .
\end{aligned}
\tag{1.9}
\]
玻色多粒子态上的配分函数由下式得到
\[
z_B(x) = x \to \frac{1}{1-x} \equiv Z_B(x) ,\qquad
z_F(x) = x \to (1+x) \equiv Z_F(x) .
\tag{1.10}
\]
而完整配分函数为 $Z = Z_B Z_F$。对玻色子，$Z_B$ 对应的是对所有「用任意多个产生算符作用所得态」求和；对费米子，$Z_F$ 对应的是对「没有费米子」与「恰有一个费米子」的态求和。若理论有多种产生算符，单字母配分函数就会有更多项。为了计算完整配分函数，我们希望把 $z_B$ 与 $z_F$ 中的每个单项式都作这一替换，然后取乘积。这由一个称为「齐次指数化」（plethystic exponentiation）的数学操作完成，即 PE。我们定义
\[
PE[f(x_i)] = \exp\left(\sum_{n=1}^{\infty}\frac{1}{n} f(x_i^n)\right),\qquad
\widetilde{PE}[f(x_i)] = \exp\left(-\sum_{n=1}^{\infty}\frac{1}{n}(-1)^n f(x_i^n)\right)
\tag{1.11}
\]
于是
\[
\boxed{\;Z_B(x) = PE[z_B(x)],\qquad
Z_F(x,\mathbf{f}) = \widetilde{PE}[z_F(x,\mathbf{f})],\qquad
Z = Z_B Z_F\;.}
\tag{1.12}
\]

\noindent\textbf{指标}\par
在含相互作用的理论中，我们将无法如此容易地算出配分函数。保持超对称的一般相互作用理论的拉氏量为
\[
L = \frac{1}{2}\dot x^2 - \frac{1}{2}W'(x)^2 + \ii\bar\psi\dot\psi - W''(x)\bar\psi\psi .
\tag{1.13}
\]
% [原文 p.5]
这里函数 $W(x)$ 称为超势。取 $W(x) = x^2/2$ 便回到超对称简谐振子。显然，对任何更高次的多项式 $W(x)$，理论都将是相互作用的。理论的谱、从而配分函数，都难以精确计算。尽管如此，仍存在一个可被精确计算的「受保护量」。
让我们注意到，所有能量 $E \neq 0$ 的态都是二重简并的：一个态是玻色的，另一个是费米子的。这纯粹是超对称的结果。设 $|\psi\rangle$ 是哈密顿量能量为 $E \neq 0$ 的一个本征态，
\[
H(\bar Q|\psi\rangle) = \bar Q H|\psi\rangle = E(\bar Q|\psi\rangle) .
\tag{1.14}
\]
除零能态之外的谱是成对的。这种非零能态的成对现象是只依赖于超对称的稳健现象，即使理论具有保持超对称的相互作用，它也应当成立。当我们打开相互作用时，理论的谱会改变，但由于超对称，各态将始终遵守成对规则。态可以变成零能，但只能成对地变；态也可以离开零能，同样只能成对地离开。这一观察促使 Witten 定义了指标 [10]，
\[
I = \Tr(-1)^F x^H .
\tag{1.15}
\]
看起来这只是在配分函数上作了一个小改动，但差别很大。$(-1)^F$ 因子使非零能处的玻色与费米贡献相互抵消，只留下零能态的贡献。对超对称谐振子，具体地
\[
I = \Tr(-1)^F x^H = 1 + (x + x^2 + \cdots) - (x + x^2 + \cdots) = 1 .
\tag{1.16}
\]
上述论证告诉我们，即使在这个理论中打开保持超对称的相互作用，指标也将始终为 1\footnote{必须确保理论的形变不会导致非归一化的真空态，因为在此类跃迁中，新的零能态可能来自无穷远。关于超对称量子力学语境下的更多讨论，见 [11]。}。换言之，相互作用理论的指标可以通过把相互作用强度设为零来计算，因为指标在光滑形变下不变。这样，前面发展的齐次指数技术就有助于计算指标。
指标只接收 $H = 0$ 的态的贡献；又因为自由理论中能量是可加的，指标也可以从「单字母指标」出发来构造。为计算自由理论的配分函数，我们曾对玻色子与费米子使用两个不同的齐次指数公式。还有一种更简单的办法，从单粒子指标构造多粒子指标：
\[
\boxed{\;I = PE[i(x)],\qquad i(x) = i_B(x) - i_F(x)\;.}
\tag{1.17}
\]
% [Tutor Note: 术语：指标（index）带 $(-1)^F$ 分次、只数零能态；配分函数（partition function）数所有态。这是全篇最易混的一对概念。]
% [Tutor Note: 易错：$(-1)^F$ 中 $F$ 是费米子数算符，玻色态取 +1、费米态取 −1；非零能态成对抵消的正是这一符号。]
% [Tutor Note: 动机：$Q$ 与 $\bar Q$ 把玻色态与费米态互相映射，(1.14) 说明 $E\neq0$ 时 $\bar Q|\psi\rangle$ 必是能量相同的伙伴态。]
% [Tutor Note: 易错：$\Z_2$ 只保留 $F \bmod 2$，因此 $(-1)^F$ 良定义；整数费米子数本身在超对称代数中并不守恒。]
% [Tutor Note: 易错：(1.8) 是普通配分函数、没有 $(-1)^F$；(1.15) 才是指标。自由谐振子两者都可算，但只有后者受保护。]
% [Tutor Note: 术语：逸度 $x = \ee^{-\beta}$，其中 $\beta$ 是 Euclidean 时间圆周长（逆温度）；$x^H$ 即 $\ee^{-\beta H}$。]
% [Tutor Note: 动机：单字母指标 $i(x) = i_B(x) - i_F(x)$ 带相对负号，正是分次的体现；PE 把它组装成多粒子指标。]
% [Tutor Note: 动机：$\widetilde{PE}$ 与 $PE$ 每项差 $(-1)^n$，恰好生成费米子的 $(1+x)$ 展开，从而 (1.10) 成立。]
% [Tutor Note: 文献：Witten 1982 [10] 提出 Witten 指标；脚注 [11] 讨论形变可能导致非归一化真空，需另行排除。]
% [Tutor Note: 易错：(1.12) 中 $\mathbf{f}$ 为原文引入的费米子逸度占位符，但 (1.9)–(1.10) 的 $z_F$ 只依赖 $x$，后文亦未再用；此处照原文保留。]
% [Tutor Note: 术语：superpotential 译「超势」；$W'$ 给出势能、$W''$ 给出费米子质量项，(1.13) 是 $\sN=1$ 量子力学的标准形式。]
% [Tutor Note: 动机：自由理论中能量可加，故多粒子谱是单粒子谱的组合；齐次指数正是这个组合运算的生成函数。]
% [Tutor Note: 语病：原文 p.4 「compute the partition function it so easily」中的 it 疑为衍字，应为「compute the partition function so easily」；照译并登记 ERRATA。]
% [Tutor Note: 语病：原文 p.5 「It is seems like」系动词重复，应为「It seems like」；已登记 ERRATA。]

% sec1b.tex — c03｜原文 p6–7｜§1.2 Global symmetry
% 编号公式 (1.18)–(1.30) 共 13 条；编号环境：习题 1.1。
% [原文 p.6]
这里，
\[
\begin{aligned}
i_B(x) &= \Tr_{(\text{bosonic letters with } H=0)} x^H ,\\
i_F(x) &= \Tr_{(\text{fermionic letters with } H=0)} x^H .
\end{aligned}
\tag{1.18}
\]
\begin{exercise}[1.1]
证明这一点。
\end{exercise}
对超对称谐振子，玻色单字母指标 $i_B(x) = 0$，费米单字母指标 $i_F(x) = 0$。因此 $i(x) = 0$。这给出完整的指标为
\[
I(x) = PE[i(x)] = 1 .
\tag{1.19}
\]
重复一遍：对相互作用的超对称理论，配分函数不易计算，特别地它是 $x$ 的函数；而指标容易计算（把耦合设为零即可），并且不依赖于 $x$（因为只有 $H=0$ 的态有贡献）。

\subsection{整体对称性}

考虑上述拉氏量的 $2N$ 个副本
\[
L = \frac{1}{2}|\dot x_i|^2 - \frac{1}{2}|x_i|^2 + \ii\bar\psi_i\dot\psi_i - \bar\psi_i\psi_i .
\tag{1.20}
\]
它具有 $\SO(2N)$ 对称性。有一种引入保持该对称性的相互作用的方法：
\[
L = \frac{1}{2}\dot x_i^2 - \frac{1}{2}\left|\frac{\partial W(x_i)}{\partial x_i}\right|^2 + \ii\bar\psi_i\dot\psi_i - \frac{\partial^2 W(x)}{\partial x_i\partial x_j}\bar\psi_i\psi_j .
\tag{1.21}
\]
取 $W = |x_i|^2$ 便得到自由理论。眼下让我们专注于自由理论。由于理论是自由的，配分函数很容易计算。它包含 $2N$ 个玻色产生算符 $a_{B,i}^\dagger$ 与 $2N$ 个费米产生算符 $a_{F,i}^\dagger$。这意味着
\[
z_B(x) = 2N x ,\qquad z_F(x) = 2N x .
\tag{1.22}
\]
利用齐次指数公式，
\[
Z(x) = \left(\frac{1+x}{1-x}\right)^{2N} .
\tag{1.23}
\]
有时记录态在整体对称性下的表示是有用的。在此情形下，我们具有 $\SO(2N)$ 对称性。作法是引入相对于 Cartan 生成元的逸度。
\[
Z(x,a_i) = \Tr x^H a_i^{J_i} ,\qquad i = 1,\ldots,N .
\tag{1.24}
\]
% [原文 p.7]
最好用单字母配分函数来计算它。由于玻色子处在 $\SO(N)$ 的基础表示中。它们的单字母配分函数为
\[
z_B(x,a_i) = x\left(a_1 + \frac{1}{a_1} + \ldots + a_n + \frac{1}{a_n}\right) = x\,\chi_{\fund}(a_i) .
\tag{1.25}
\]
这里 $\chi_{\fund}(a_i)$ 是基础表示的特征标。一般地，一个表示的特征标定义为
\[
\chi_R(a_i) = \sum_{\rho\in R}\prod_i a_i^{\rho_i} \equiv \sum_{\rho\in R} a^\rho ,
\tag{1.26}
\]
其中 $\rho = (\rho_1,\ldots,\rho_N)$ 是权向量，求和遍历该表示的全部权向量。有时我们会把 $\prod_i a_i^{\rho_i} = a^\rho$ 缩写。回忆一下，伴随表示的非零权向量称为根。于是我们有，
\[
\begin{aligned}
z_B(x,a_i) &= x\left(a_1 + \frac{1}{a_1} + \ldots + a_n + \frac{1}{a_n}\right) = x\,\chi_{\fund}(a_i)\\
Z_B(x,a_i) &= PE[z(x,a)] = \prod_{i=1}^{N}\frac{1}{(1-x a_i)(1-x/a_i)}\\
z_F(x,a_i) &= x\left(a_1 + \frac{1}{a_1} + \ldots + a_n + \frac{1}{a_n}\right) = x\,\chi_{\fund}(a_i)\\
Z_F(x,a_i) &= \widetilde{PE}[z(x,a)] = \prod_{i=1}^{N}(1+x a_i)(1+x/a_i)\\
Z(x,a_i) &= Z_B(x,a_i)Z_F(x,a_i) = \prod_{i=1}^{N}\frac{(1+x a_i)(1+x/a_i)}{(1-x a_i)(1-x/a_i)} \equiv \prod_{\rho}\frac{(1+x a^\rho)}{(1-x a^\rho)} .
\end{aligned}
\tag{1.27}
\]
显然，当取 $a_i = 1$ 时，$Z(x,a_i)$ 退化为未精细化的配分函数 $Z(x)$。记录相对于 Cartan 生成元的荷的好处在于，我们可以计算具有给定整体对称性表示的态上的配分函数。这要用到特征标的正交性，
\[
\frac{1}{|W|}\oint\prod_{i=1}^{N}\frac{\dd a_i}{2\pi\ii a_i}\,\Delta(a_i)\chi_R(a_i)\chi_{R'}(a_i) = \delta_{RR'} .
\tag{1.28}
\]
其中 $N$ 是群的秩，$|W|$ 是相伴 Weyl 群的阶，$\Delta(a_i)$ 称为 Vandermonde 行列式。Vandermonde 行列式有下述用齐次指数表示的普遍公式\footnote{需要正则化才能使齐次指数良定义。}
\[
\Delta(a_i) = \left(PE\Big[\sum_{\alpha} a^\alpha\Big]\right)^{-1} = \prod_{\alpha}(1-a^\alpha) ,
\tag{1.29}
\]
% [原文 p.8]
其中 $\alpha = (\alpha_1,\ldots,\alpha_N)$ 是李代数的根。现在很清楚如何投影到具有给定表示的态上，
\[
Z(x,a_i)\big|_{R} = \frac{1}{|W|}\oint\prod_{i=1}^{N}\frac{\dd a_i}{2\pi\ii a_i}\,\Delta(a_i)Z(x,a_i)\chi_R(a_i) .
\tag{1.30}
\]
在计算规范理论的指标时，我们将需要投影到具有平庸表示的态上。平庸表示的特征标为 $\chi_R(a_i) = 1$。
% [Tutor Note: 术语：$i_B$、$i_F$ 是「单字母指标」，与 $z_B$、$z_F$（单字母配分函数）只差 $(-1)^F$ 分级；习题 1.1 要证的正是 (1.17) 在本例中给出 $I=1$。]
% [Tutor Note: 动机：$2N$ 个副本即 $N$ 个复谐振子，$\SO(2N)$ 来自把这 $2N$ 个实分量整体旋转。]
% [Tutor Note: 易错：(1.20) 用复变量 $x_i$ 写成 $|x_i|^2$，所以「$2N$ 个副本」指 $2N$ 个实自由度，不是 $2N$ 个复变量。]
% [Tutor Note: 易错：$\chi_R$ 是群论「特征标」，不要与 $\sI$（指标 index）混淆；本译本统一译作「特征标」。]
% [Tutor Note: 术语：$a^\rho$ 只是 $\prod_i a_i^{\rho_i}$ 的缩写，$\rho$ 是权向量而非单个整数。]
% [Tutor Note: 动机：加入 Cartan 逸度后 $Z(x,a_i)$ 的展开按 $\SO(2N)$ 表示分解；令全部 $a_i=1$ 即退回未精细化配分函数。]
% [Tutor Note: 易错：原文此处写 $\SO(N)$（前文为 $\SO(2N)$），而 (1.25) 的展开恰是 $\SO(2N)$ 矢量表示的特征标；疑为原文笔误，已登记 ERRATA，译文照录。]
% [Tutor Note: 文献：特征标正交性 (1.28) 是 Weyl 积分公式的投影形式；$\frac{1}{|W|}$ 与 Haar 测度配套出现。]
% [Tutor Note: 术语：$\Delta(a_i)=\prod_\alpha(1-a^\alpha)$ 遍历全部根（含正负根），故等于 Weyl 分母的绝对值平方；脚注 2 说明 PE 求逆需先正则化。]
% [Tutor Note: 易错：(1.29) 对 $\alpha$ 的乘积遍历全部根，不能只取正根，否则会漏掉 $1-a^{-\alpha}$ 因子。]


\section{超共形指标}
% sec2a.tex — c04｜原文 p8–12｜§2 超共形指标（开头含 §1 末尾的习题 1.2 与两段收尾）
% 编号公式 (2.1)–(2.16) 共 16 条；编号环境：习题 1.2、2.1；表 1（复刻为 tabular）。
% 结构说明：分块表把 p8 的习题 1.2 归入本文件；原文中它在 §2 标题之前，故编译后会落在
% §2 标题之下。此处按分块表施工，已在块末批注与回报中列为未决。
% [原文 p.8]
\begin{exercise}[1.2]
在 $x$ 展开中以及用 Cauchy 定理对 $\SO(2)$ 计算这个积分。
\end{exercise}

再次，单字母指标为 $i_B(x,a)=0$ 且 $i_F(x,a)=0$。这意味着 $I(x,a)=PE[i_B(x)-i_F(x)]=1$。注意，在指标中只能开启那些与超荷对易的荷的逸度。这是因为指标的稳健性依赖于非零能态上的玻色—费米抵消。若开启与 $Q$ 不对易的荷的逸度，抵消将不成立，指标在量子修正下也不会受到保护。

到目前为止，我们还没有遇到具有非平庸指标的理论，但我们已经发展出在指标非平庸时有助于计算它的工具。现在我们将把注意力转向超共形指标。

四维超共形场论的指标定义为该理论在径向量子化下的 Witten 指标。取 $Q\equiv Q_-$ 为其一 Poincaré 超荷，$Q^\dagger=S_+$ 为共轭的共形超荷。概略地说，指标定义为 [12–14]
\[
\sI(\mu_i)=\Tr(-1)^F x^{\delta}\mu_i^{\mathcal{M}_i} ,
\tag{2.1}
\]
其中迹取在该理论于 $S^3$ 上的 Hilbert 空间上，$\delta\equiv\frac12\{Q,Q^\dagger\}$。由于 $\delta>0$ 的态以玻色/费米对的形式出现，只有 $\delta=0$ 的态有贡献，且指标与 $x$ 无关。如前所述，费米—玻色抵消只有在额外的逸度 $\mu_i$ 仅对与 $Q$ 对易的那些荷 $\mathcal{M}_i$ 开启时才成立。与量子力学的情形不同，该理论有无穷多个 $\delta=0$ 的态。不仅如此，正如我们很快将看到的，该理论还有无穷多个 $\delta=0$ 的「单字母」。逸度 $\mu_i$ 的引入既是为了正则化这一发散，也是为了实现更精细的计数。

对 $\sN=1$，超荷为 $\{Q_\alpha, S^\alpha\equiv Q^{\dagger\alpha}, \tilde Q_{\dot\alpha}, \tilde S^{\dot\alpha}\equiv \tilde Q^{\dagger\dot\alpha}\}$，其中 $\alpha=\pm$ 与 $\dot\alpha=\dot\pm$ 分别是 $\SU(2)_1$ 与 $\SU(2)_2$ 的指标，而 $\SU(2)_1\times \SU(2)_2=\mathrm{Spin}(4)$ 是 $S^3$ 的等距群。相关的反对易关系为
% [原文 p.9]
\[
\{Q_\alpha, Q^{\dagger\beta}\}=\Delta+2M_\alpha^\beta+\frac32 r ,
\tag{2.2}
\]
\[
\{\tilde Q_{\dot\alpha}, \tilde Q^{\dagger\dot\beta}\}=\Delta+2\tilde M_{\dot\alpha}^{\dot\beta}-\frac32 r ,
\tag{2.3}
\]
其中 $\Delta$ 是共形维数，$M_\alpha^\beta$ 与 $\tilde M_{\dot\alpha}^{\dot\beta}$ 是 $\SU(2)_1$ 与 $\SU(2)_2$ 生成元，$r$ 是 $\U(1)_r$ R 对称性的生成元。在我们的约定中，$Q$ 的 $r=-1$，而 $\tilde Q$ 的 $r=+1$；当然，dagger 操作会翻转 $r$ 的符号。

取超荷 $Q\equiv Q_-$ 来定义指标，我们得到 $\delta=\Delta-2j_1+\frac32 r$。$Q$ 在超共形代数 $\SU(2,2|\sN)$ 中的对易子为 $\SU(2,1|\sN-1)$。对 $\sN=1$，对易子为 $\SU(2,1)$。它的秩为 2。我们可以在这个空间中选择任意一组荷作为 $\mathcal{M}_i$ 来精细化指标。为了与已有文献一致，我们取其为 $\frac13(\Delta+j_1)\pm j_2$。于是指标定义为
\[
\begin{aligned}
\sI(p,q)&\equiv\Tr(-1)^F p^{\frac13(\Delta+j_1)+j_2}q^{\frac13(\Delta+j_1)-j_2}\\
&=\Tr(-1)^F p^{j_1-\frac12 r+j_2}q^{j_1-\frac12 r-j_2},\qquad \delta=\Delta-2j_1+\frac32 r ,
\end{aligned}
\tag{2.4}
\]
在第二个等号中，我们利用只有 $\delta=0$ 的态对指标有贡献这一事实，把荷 $\frac13(\Delta+j_1)$ 换成了 $j_1-\frac12 r$。这样做是因为 $r$ 荷比共形维数更容易确定。

\subsection{手征多重态}

我们从单个自由手征多重态 $\Phi$ 的理论开始。它满足条件 $\bar Q_{\dot\alpha}\Phi=0$。手征多重态的超空间展开为
\[
\Phi=\phi+\sqrt2\,\theta\psi+\ii\,\theta^\dagger\bar\sigma^\mu\theta\,\partial_\mu\phi .
\tag{2.5}
\]
复共轭是一个反手征场，即 $Q_\alpha\bar\Phi=0$。$\phi$ 的 $r$ 荷为 $2/3$（你知道为什么吗？）。

在径向量子化中，理论的 Hilbert 空间同构于局域算符的谱。这些算符由手征多重态中的场及其导数构造。这意味着，量子力学中产生算符的对应物是场及其导数。它们是
\[
\phi,\ \partial_\mu\phi,\ \partial_\mu\partial_\nu\phi,\ \ldots,\qquad
\psi,\ \partial_\mu\psi,\ \partial_\mu\partial_\nu\psi,\ \ldots
\tag{2.6}
\]
及其共轭。这里导数以所有可能的缩并作用，但带有来自运动方程的约束 $\partial_\mu\partial^\mu\phi=0$。类似地，费米子也满足运动方程 $\partial_{\alpha\dot\alpha}\psi^\alpha=0$。有了这些约束，我们在表 1 中列出字母并枚举「单字母指标」。由此，我们可以容易地计算单字母指标，
\[
i_\phi(p,q)=\frac{(pq)^{\frac13}-(pq)^{\frac23}}{(1-p)(1-q)} .
\tag{2.7}
\]
% [原文 p.10]
\noindent\textbf{表 1.} $\sN=1$ 手征多重态的「字母」及其对指标的贡献。其中 $\delta=\Delta-2j_1+\frac32 r$。先验地，我们必须考虑自由运动方程 $\partial\psi=0$ 与 $\Box\phi=0$，它们对可能的词施加约束；但我们看到，在此情形下运动方程的 $\delta\neq0$，因此它们对单字母指标没有贡献。最后有两个时空导数对指标有贡献（即 $\delta=0$），它们在场上形成的塔负责单字母指标的分母。

\begin{center}
\begin{tabular}{C{4.6em} C{1.4em} C{1.4em} C{1.4em} C{1.4em} C{1.4em} C{2.2em}}
\toprule
字母 & $\Delta$ & $j_1$ & $j_2$ & $r$ & $\delta$ & $\sI$ \\
\midrule
$\phi$ & $1$ & $0$ & $0$ & $\frac23$ & $2$ & \\
$\psi_+$ & $\frac32$ & $\frac12$ & $0$ & $-\frac13$ & $0$ & $-(pq)^{\frac23}$ \\
$\psi_-$ & $\frac32$ & $-\frac12$ & $0$ & $-\frac13$ & $2$ & \\
$\partial^\alpha_{\dot+}\psi_\alpha$ & $\frac52$ & $0$ & $\frac12$ & $-\frac13$ & $2$ & \\
$\partial^\alpha_{\dot-}\psi_\alpha$ & $\frac52$ & $0$ & $-\frac12$ & $-\frac13$ & $2$ & \\
$\Box\phi$ & $3$ & $0$ & $0$ & $\frac23$ & $4$ & \\
\midrule
$\bar\phi$ & $1$ & $0$ & $0$ & $-\frac23$ & $0$ & $(pq)^{\frac13}$ \\
$\bar\psi_{\dot+}$ & $\frac32$ & $0$ & $\frac12$ & $\frac13$ & $2$ & \\
$\bar\psi_{\dot-}$ & $\frac32$ & $0$ & $-\frac12$ & $\frac13$ & $2$ & \\
$\partial^{\dot\alpha}_{+}\bar\psi_{\dot\alpha}$ & $\frac52$ & $\frac12$ & $0$ & $\frac13$ & $2$ & \\
$\partial^{\dot\alpha}_{-}\bar\psi_{\dot\alpha}$ & $\frac52$ & $-\frac12$ & $0$ & $\frac13$ & $4$ & \\
$\Box\bar\phi$ & $3$ & $0$ & $0$ & $-\frac23$ & $2$ & \\
\midrule
$\partial_{+\dot+}$ & $1$ & $\frac12$ & $\frac12$ & $0$ & $0$ & $p$ \\
$\partial_{-\dot+}$ & $1$ & $-\frac12$ & $\frac12$ & $0$ & $0$ & $q$ \\
$\partial_{+\dot-}$ & $1$ & $\frac12$ & $-\frac12$ & $0$ & $2$ & \\
$\partial_{-\dot-}$ & $1$ & $-\frac12$ & $-\frac12$ & $0$ & $2$ & \\
\bottomrule
\end{tabular}
\end{center}

分母编码了两个 $\delta=0$ 的时空导数的作用。如果手征多重态在某个整体对称性的表示 $R$ 中变换，那么我们可以用 Cartan 生成元的逸度来精细化单字母指标。注意，如果手征多重态在表示 $R$ 中变换，则反手征多重态在共轭表示 $\bar R$ 中变换。
\[
i_\phi(a_i;p,q)=\frac{(pq)^{\frac13}\chi_{\bar R}(a_i)-(pq)^{\frac23}\chi_R(a_i)}{(1-p)(1-q)} .
\tag{2.8}
\]
\begin{exercise}[2.1]
相对于 $\tilde Q_{\dot+}$ 定义指标，并仿照类似的表格计算 $\Phi$ 的单字母指标。
\end{exercise}
全指标通过齐次指数化得到。对在 $\U(1)$ 对称性下荷为 1 的单个多重态
% [原文 p.11]
，其指标为
\[
\sI_\phi(a;p,q)=\Gamma\big((pq)^{\frac13}a^{-1};p,q\big),\qquad
\Gamma(z;p,q)\equiv\prod_{i,j=1}^{\infty}\frac{1-p^{i+1}q^{j+1}/z}{1-p^iq^jz} .
\tag{2.9}
\]
特殊函数 $\Gamma(z;p,q)$ 称为椭圆 gamma 函数 [15]，在数学文献中作为通常 Gamma 函数的「椭圆」推广而被研究。椭圆 gamma 函数在超共形指标理论中起重要作用，因为手征多重态是四维任何超共形拉氏量不可或缺的组成部分。我们几乎总是略去椭圆 Gamma 函数的自变量 $(p,q)$，并记 $\Gamma(z;p,q)\to\Gamma(z)$。如果手征多重态在表示 $R$ 中变换，则相应的指标为
\[
\sI_\phi(a_i;p,q)=\prod_{\rho\in R}\Gamma\big((pq)^{\frac13}a^\rho\big) .
\tag{2.10}
\]
如果手征多重态参与超势相互作用，那么它的 R 荷可以不同于 $2/3$。它由要求超势具有 R 荷 2 来确定。如果理论具有阿贝尔整体对称性，这一约束不能完全固定 R 荷，因为这些对称性可以与这个假定的 R 荷混合。在那种情形下，R 荷由要求 R 对称性与任何阿贝尔整体对称性都没有混合反常来确定。这一程序称为「a-极大化」[16,17]。荷为 $r$ 的手征多重态的单字母指标，是遍及自由情形下对指标有贡献的字母求和得到的；但这些字母按其 R 荷对指标有贡献。这称为 Romelsberger 处方 [13,14]。这一奇特处方后来通过在 $S^3\times S^1$ 上沿 RG 流保持超对称而被理解 [18]。

由表 1，手征多重态中对指标有贡献的字母是 $\psi_+$ 与 $\bar\phi$。若手征多重态的 R 荷为 $r$，则这两个字母的 R 荷分别为 $r-1$ 与 $-r$。相应地，它们将分别贡献 $-(pq)^{1-\frac12 r}$ 与 $(pq)^{\frac12 r}$。这给出 R 荷为 $r$ 的手征多重态的指标为
\[
i_{\phi_r}(p,q)=\frac{(pq)^{\frac12 r}-(pq)^{1-\frac12 r}}{(1-p)(1-q)} .
\tag{2.11}
\]
多粒子指标为
\[
\sI_{\phi_r}(a;p,q)=\Gamma\big((pq)^{\frac12 r}a\big) .
\tag{2.12}
\]

\paragraph{例子} 有了这些工具，我们已经可以计算一些只涉及手征多重态的物理理论的指标。本小节的例子取自综述 [8]。

\paragraph{有质量理论：} 考虑带质量项的单个手征多重态理论，即 $W(\Phi)=m\Phi^2$。这把手征场的 R 荷固定为 1。
\[
\sI_{\phi_1}(p,q)=\Gamma\big((pq)^{\frac12}\big)=1 .
\tag{2.13}
\]
% [原文 p.12]
这与物理预期一致：有质量理论应当只有一个超对称基态。

作为轻微的推广，可以考虑两个手征多重态的理论，其质量型超势为 $W(\Phi_i)=m\Phi_1\Phi_2$。这个超势不足以完全固定 $\Phi_{1,2}$ 的 R 荷，但它确实告诉我们 $r_1+r_2=2$。仅凭这一约束
\[
\sI(a;p,q)=\Gamma\big((pq)^{\frac12 r}a\big)\Gamma\big((pq)^{1-\frac12 r}a^{-1}\big)=1 .
\tag{2.14}
\]
这也在预料之中，因为该超势使两个手征多重态都获得能隙，从分量形式写出的拉氏量可以清楚地看出这一点。

\paragraph{F 项超对称破缺} 考虑一个简单的 O'Raifeartaigh 型模型，其超势为线性 $W(\Phi)=\eta\Phi$。该理论自发破缺超对称，因为该超势给基态以非零能量。让我们看看这在指标中如何体现。该超势把 $\Phi$ 的 R 荷固定为 2，且不允许任何整体对称性。
\[
\sI_{\phi_2}(p,q)=\Gamma(pq)=0 .
\tag{2.15}
\]
指标中的零来自 $\psi_+$ 的贡献变为 $-(pq)^{1-\frac12 r}\to-1$。这意味着带 $\psi_+$ 的态 $\psi_+|0\rangle$ 与 $|0\rangle$ 具有相同的能量。换言之，基态以玻色—费米对的形式出现，因此不受保护。这也意味着 $\psi_+$ 应被解释为超对称自发破缺的 Goldstino 模。任何自发破缺超对称的 O'Raifeartaigh 型模型总有一个荷为 2 的中性手征多重态。这导致 $\sI=0$。

\paragraph{非平庸手征环} 具有非平庸手征环的最简单模型是单个手征多重态的理论，其超势为 $W(\Phi)=\Phi^{n+1}$。手征环关系为 $W'(\Phi)=0=\Phi^n$。超势意味着场 $\Phi$ 具有 R 荷 $2/(n+1)$，并且它关于一个离散对称性 $\mathbb{Z}_{n+1}$ 带荷 1。我们可以开启味对称性的逸度 $g$ 来计算指标。由于味对称性只是离散的 $\mathbb{Z}_{n+1}$，逸度满足 $g^{n+1}=1$。单字母指标为
\[
i(g;p,q)=\frac{(pq)^{\frac{1}{n+1}}g-(pq)^{1-\frac{1}{n+1}}g^{-1}}{(1-p)(1-q)}
=\frac{(pq)^{\frac{1}{n+1}}g-\big((pq)^{\frac{1}{n+1}}g\big)^n}{(1-p)(1-q)} .
\tag{2.16}
\]
来自费米字母 $\psi_+$ 的贡献恰好是玻色字母贡献的 $n$ 倍，且符号相反。这与手征环关系 $\Phi^n=0$ 相符。

% [Tutor Note: 术语：本节的「指标」$\sI$ 即四维 SCFT 的 Witten 指标在径向量子化下的版本；与 §1 的配分函数 $Z(x)$ 不同，它能对超共形代数做精细计数。]
% [Tutor Note: 动机：$\delta\equiv\frac12\{Q,Q^\dagger\}$ 扮演「能量」；$\delta>0$ 的态成对抵消，故指标只数 BPS 态，且与 $x$ 无关。]
% [Tutor Note: 易错：超共形指标中 $x^\delta$ 的 $\delta$ 已换成 $\Delta-2j_1+\frac32 r$，不要与 §1 的哈密顿量 $H$ 混同。]
% [Tutor Note: 易错：$\mu_i$ 只能对与 $Q$ 对易的荷 $\mathcal{M}_i$ 开启，否则玻色—费米抵消失效，指标不再受保护。]
% [Tutor Note: 术语：$S^\alpha\equiv Q^{\dagger\alpha}$ 是共形超荷；$\alpha$ 与 $\dot\alpha$ 分别属于 $\SU(2)_1$、$\SU(2)_2$，合成 $\mathrm{Spin}(4)$，不要与 R 对称性混淆。]
% [Tutor Note: 动机：$Q$ 的对易子 $\SU(2,1|\sN-1)$ 的秩决定可开启的逸度个数；$\sN=1$ 时秩 2，故 $\sI$ 依赖 $p,q$。]
% [Tutor Note: 易错：(2.4) 第二个等号只在 $\delta=0$ 的态上成立，它把较难确定的共形维数 $\Delta$ 换成 $r$ 荷。]
% [Tutor Note: 术语：$\Gamma(z;p,q)$ 是椭圆 gamma 函数，不要与普通 gamma 函数 $\Gamma(z)$ 混淆；后者是 $p,q\to0$ 的极限。]
% [Tutor Note: 动机：表 1 的分母 $(1-p)(1-q)$ 来自两个 $\delta=0$ 的时空导数 $\partial_{+\dot+}$、$\partial_{+\dot-}$ 的塔。]
% [Tutor Note: 易错：表 1 第三组第 2、3 行的 $\delta$ 与 $\sI$ 列疑有笔误：按 $\delta=\Delta-2j_1+\frac32 r$，$\partial_{-\dot+}$ 的 $\delta$ 应为 $2$、$\partial_{+\dot-}$ 的 $\delta$ 应为 $0$（$q$ 应随 $\delta=0$ 行）；译文照录，已登记 ERRATA。]
% [Tutor Note: 易错：表 1 只列「单字母」；词的贡献要再把单字母指标做齐次指数化（(2.9) 的 $\Gamma$ 就是这一塔的求和）。]
% [Tutor Note: 易错：自由手征多重态 $r=2/3$；有超势时 $r$ 由 $W$ 的 R 荷 2 与 a-极大化决定，不能一律取 $2/3$。]
% [Tutor Note: 术语：Romelsberger 处方＝用自由理论的字母、但按相互作用理论的 R 荷赋权；文献 [13,14] 提出，[18] 用 $S^3\times S^1$ 上的超对称流解释。]
% [Tutor Note: 易错：$\sI=0$ 不是「没有态」，而是基态成玻色—费米对、超对称自发破缺；此时 $\psi_+$ 是 Goldstino。]
% [Tutor Note: 易错：$\mathbb{Z}_{n+1}$ 是离散味对称性，逸度满足 $g^{n+1}=1$，不能当作连续 $\U(1)$ 的逸度使用。]
% [Tutor Note: 文献：[15] 椭圆 gamma 函数；[16,17] a-极大化；[8] 综述中的例子。]
% [Tutor Note: 结构提醒：习题 1.2 原文在 §1 末、§2 标题之前，但按分块表归入本文件，编译后会出现在 §2 标题之下。]
% [Tutor Note: 语病：原文 p.11 「the index it studied」缺系动词，应为「the index it is studied」或「the index studied」；照译并登记 ERRATA。]
% [Tutor Note: 语病：原文 p.12 「the this superpotential」冠词重复，应为「this superpotential」；已登记 ERRATA。]

% sec2b.tex — c05｜原文 p13–18（含 p19 顶端一句收尾）｜§2.2 规范理论的指标、§2.3 重组不变性、§2.4 拉氏量概览
% 编号公式 (2.17)–(2.29) 共 13 条；编号环境：习题 2.2、2.3；表 2（tabular 复刻）。
% 说明：§2.4 末句跨页到 p19 顶端，为保证句子完整，本块含 p19 开头一句；sec3a 从 §3 起。
% [原文 p.13]
\subsection{规范理论的指标}

除了 $\sN=1$ 手征多重态之外，四维超共形理论中随处可见的另一个多重态是 $\sN=1$ 矢量多重态。矢量多重态的指标可以用同样的方式计算，即先计算单字母指标，再取齐次指数。我们在表 2 中列出了 $\delta=0$ 的矢量多重态字母。
\begin{exercise}[2.2]
相对于 $\tilde Q_{\dot+}$ 计算矢量多重态的单字母指标。
\end{exercise}
\noindent\textbf{表 2.} 本表只列出 $\delta=0$ 的字母。这些是对指标有贡献的字母。

\begin{center}
\begin{tabular}{C{4.6em} C{1.4em} C{1.4em} C{1.4em} C{1.4em} C{1.4em} C{2.2em}}
\toprule
字母 & $\Delta$ & $j_1$ & $j_2$ & $r$ & $\delta$ & $\sI$ \\
\midrule
$\bar\lambda_{\dot+}$ & $\frac32$ & $0$ & $\frac12$ & $-1$ & $0$ & $-p$ \\
$\bar\lambda_{\dot-}$ & $\frac32$ & $0$ & $-\frac12$ & $-1$ & $0$ & $-q$ \\
$F_{++}$ & $2$ & $1$ & $0$ & $0$ & $0$ & $pq$ \\
$\partial^{\dot\alpha}_{+}\bar\lambda_{\dot\alpha}$ & $\frac52$ & $\frac12$ & $0$ & $-1$ & $0$ & $pq$ \\
$\partial_{+\dot+}$ & $1$ & $\frac12$ & $\frac12$ & $0$ & $0$ & $p$ \\
$\partial_{-\dot+}$ & $1$ & $-\frac12$ & $\frac12$ & $0$ & $0$ & $q$ \\
\bottomrule
\end{tabular}
\end{center}

因此，矢量多重态的单字母指标为
\[
i_V(a_i:p,q)=\frac{-p-q+2pq}{(1-p)(1-q)}\chi_{\adj}(a_i)
=\Big(-\frac{p}{1-p}-\frac{q}{1-q}\Big)\Big(\sum_\alpha a^\alpha+N\Big) .
\tag{2.17}
\]
这里我们已把单字母指标乘以规范群伴随表示的特征标。这是因为矢量多重态总处于伴随表示。伴随表示的特征标为 $\chi_{\adj}(a_i)=\sum_\alpha a^\alpha+N$，其中 $\alpha$ 是根向量。加性因子 $N$ 对应权向量为零的 Cartan 生成元。

既然已经知道手征多重态与矢量多重态的指标，我们就做好准备计算规范理论的指标了。依赖于物质内容，低能下的超对称规范理论流向非平庸的超共形场论。很可能无法精确计算它们的谱。然而超共形指标容易计算。诀窍在于意识到指标在规范耦合改变下不变。强耦合规范理论的指标可以通过把规范耦合设为零来计算。

% [原文 p.14]
这意味着我们只想枚举由 $\delta=0$ 的字母构造出的词。此类词上的配分函数就是单字母指标的齐次指数。这些就是四维中能构造出的全部超对称局域算符\footnote{在三维修与五维中，还有额外局域算符不能用拉氏量中基本场的字母构造。在三维它们称为「单极子算符」，在五维称为「瞬子算符」。计入它们意味着分别对 $S^2$ 与 $S^4$ 上非平凡磁通的分区求和。本讲义只关注四维理论。}。然而在规范理论中，我们只对构造规范变换下不变的算符感兴趣。这些算符是单态，即在规范群的平庸表示下变换。因此，规范理论的指标为
\[
\sI(b_k;p,q)=\frac{1}{|W|}\oint\Big(\prod_{i=1}^{N}\frac{\dd a_i}{2\pi\ii a_i}\Big)\Delta(a_i)\,\sI_V(a_i;p,q)\prod_{\phi_i}\sI_{\phi_i}(a_i,b_k;p,q) .
\tag{2.18}
\]
这里 $a_i$ 是规范对称性的逸度，$b_k$ 是手征多重态可能带有的所有其它对称性（整体对称性）的逸度。

利用 $\Delta(a_i)=(PE[\sum_\alpha a^\alpha])^{-1}$ 这一事实，我们可以化简组合
\[
\begin{aligned}
\Delta(a_i)\sI_V(a_i;p,q)&=PE\Big[\Big(-\frac{p}{1-p}-\frac{q}{1-q}\Big)\chi_{\adj}(a_i)-\sum_\alpha a^\alpha\Big]\\
&=PE\Big[\frac{pq-1}{(1-p)(1-q)}\Big(\sum_\alpha a^\alpha\Big)+N\Big(-\frac{p}{1-p}-\frac{q}{1-q}\Big)\Big]\\
&=\kappa^N\prod_\alpha\Gamma(pq\,a^\alpha),\qquad \kappa\equiv(p,p)(q,q),\quad (a,q)\equiv\prod_{i=0}^{\infty}(1-aq^i) .
\end{aligned}
\tag{2.19}
\]
函数 $(a,q)$ 称为 Pochhammer 符号。我们用简写 $\kappa$ 表示单个 $\U(1)$ 矢量多重态的指标。

利用这一化简，规范理论的指标可以完全用椭圆 gamma 函数（以及 Pochhammer 符号）表示出来。
\[
\begin{aligned}
\sI(b_k;p,q)&=\frac{\kappa^N}{|W|}\\
&\quad\times\oint\frac{\dd a_i}{(2\pi\ii)^N}\Gamma(pq\,a^\alpha)\prod_{\phi_i}\prod_{\rho\in R_i^G,\rho'\in R'^F_i}\Gamma\big((pq)^{\frac12 r_i}a^\rho b^{\rho'}\big) .
\end{aligned}
\tag{2.20}
\]
这里我们给手征多重态指定了任意的 R 荷 $r_i$。如前所述，它由超势确定，此外可能还要用 a-极大化。

\paragraph{例子} 现在我们给出一个 $\sN=1$ 规范理论中指标计算的例子，以说明超共形指标的威力。考虑 $\sN=1$、带三味基础与反基础夸克的 $\SU(2)$ 规范理论。该理论具有 $\SU(3)\times\SU(3)\times\U(1)$ 整体对称性。由于阿贝尔因子的存在，我们需要借助 a-极大化来确定正确的 R 荷。它们的结果为 $1/3$。
% [原文 p.15]
利用规范理论指标的一般公式 (2.20)，我们把该规范理论的指标写成积分
\[
\sI=\kappa\oint\frac{\dd z}{4\pi\ii z}\frac{1}{\Gamma(z^{\pm2})}\prod_{i=1}^{3}\Gamma\big((pq)^{\frac16}b\,u_i z^{\pm1}\big)\Gamma\big((pq)^{\frac16}b^{-1}v_i z^{\pm1}\big) .
\tag{2.21}
\]
这里 $\prod_{i=1}^3 u_i=\prod_{i=1}^3 v_i=1$，这些逸度参数化味对称性 $\SU(3)_u\times\SU(3)_v$——它转动基础与反基础夸克——而 $b$ 参数化重子的 $\U(1)_b$。通常，作为物理学家，我们会令 $p=xy$、$q=x/y$，把该积分按 $x$ 作 Taylor 级数展开来处理；但这个特定的积分已在数学文献中被研究过。引人注目的是，它允许闭式表达式
\[
\sI=\prod_{i<j}\Gamma\big((pq)^{\frac13}t_it_j\big) .
\tag{2.22}
\]
这里 $\{t_i\}=\{bu_i,b^{-1}v_i\}$。这称为 Spiridonov 的 beta 函数恒等式 [19]。
\begin{exercise}[2.3]
在 $x$ 的 Taylor 级数展开中验证恒等式 (2.22)。
\end{exercise}
这一形式极具启发性。它是十五个手征多重态的指标，这些多重态与一个超势相互作用，该超势赋予每个多重态 R 荷 $2/3$。此外，$u$ 与 $v$ 逸度组合成六个 $t_i$，使得指标在它们的置换下不变，这暗示该理论具有 $\SU(6)$ 对称性，其逸度为 $t_i$（对称性是 $\SU(6)$ 而非 $\U(6)$，因为 $\prod_i t_i=1$）。这样的超势正是由这十五个手征场构造出的 $6\times6$ 反对称矩阵的 Pfaffian。

这暗示规范理论与仅含手征多重态的理论之间存在对偶。这确实成立。这是众多著名的 Seiberg 对偶之一 [20]。碰巧这一特定对偶对超共形指标的推论已在数学文献中被独立研究过。但更一般地，以这种方式，所有 Seiberg 对偶都成为非平凡数学积分恒等式的来源\footnote{这些理论在大 $N$ 极限下的超共形指标已在 [21] 中研究。}。事实上，这类积分（称为椭圆超几何积分）已被数学家单独研究。参见 [22] 的导引与文献。更重要的是，从物理学家的观点看，超共形指标的匹配为非凡的非微扰对偶猜想提供了非平凡的检验。

\subsection{重组不变性}

在超对称量子力学中，Hilbert 空间构成超对称代数 $\SU(1|1)$ 的一个表示。有两类不可约表示。一个二维长多重态 $\{|\chi\rangle,|\psi\rangle\}$
\[
Q|\chi\rangle=|\psi\rangle,\qquad Q|\psi\rangle=0
\tag{2.23}
\]
% [原文 p.16]
以及一个一维短多重态 $\{|\chi\rangle\}$
\[
Q|\phi\rangle=0 .
\tag{2.24}
\]
仅从代数 $\{Q,Q^\dagger\}=H$ 就容易看出 $H|\phi\rangle=0$。若对理论作光滑的超对称形变，长多重态的能量可能降到零，在这一点它分裂成两个短多重态。反之，若在零能处有两个短多重态，则光滑的超对称形变一般会使它们重组形成一个长多重态。Witten 指标的关键性质是它在重组下不变，这正是它在光滑超对称形变下受保护的原因。

人们可以对超共形指标提出同样的表示论问题，并把它理解为一个在超共形代数所允许的一切重组下不变的对象。Lorenz Eberhardt 已经讲过超共形代数及其长、短表示。我们将为此借鉴那一讨论。虽然重组规则最早在 [23] 中讨论，我们将遵循 [12] 的叙述。让我们回忆，

三类缩短条件
\[
\begin{aligned}
A_1:\ \epsilon^{\alpha\beta}\mathcal{Q}_\alpha([j_1,j_2]^r_\Delta)_\beta=0
&\quad\Rightarrow\quad \Delta=2+2j_1-\frac32 r
&\qquad\text{零态：}\quad &[j_1-\tfrac12,j_2]^{r-1}_{\Delta+\frac12}\\
A_2:\ \epsilon^{\alpha\beta}\mathcal{Q}_\alpha\mathcal{Q}_\beta[0,j_2]^r_\Delta=0
&\quad\Rightarrow\quad \Delta=2-\frac32 r
&\qquad\text{零态：}\quad &[0,j_2]^{r-2}_{\Delta+1}\\
B_1:\ \mathcal{Q}_\alpha[0,j_2]^r_\Delta=0
&\quad\Rightarrow\quad \Delta=-\frac32 r
&\qquad\text{零态：}\quad &[\tfrac12,j_2]^{r-1}_{\Delta+\frac12}
\end{aligned}
\tag{2.25}
\]
我们采用了 [24] 中分类缩短条件的约定。考虑 $A_1$ 型缩短条件：当超共形多重态的原初态饱和幺正性界 $\Delta=2+2j_1-\frac32 r$ 时，态 $[j_1-\frac12,j_2]^{r-1}_{\Delta+\frac12}$ 变为零态。有趣的是，容易验证该零态也处在幺正性界上。如果你现在把带零态的那个多重态视为以零态为原初态的多重态，那么这个多重态也有零态，并且同样处在幺正性界上。在这一过程中 $j_1$ 的值不断减小，链条在 $j_1=0$ 处终止。在那种情形下，零态满足 $B_1$ 型缩短条件。总之，我们得到一列短多重态，使得任意两个相邻多重态都能重组形成一个长多重态，
\[
[j_1,j_2]^r_\Delta\to[j_1-\tfrac12,j_2]^{r-1}_{\Delta+\frac12}\to\ldots\to[0,j_2]^{r-2j_1}_{\Delta+j_1}\to[0,j_2]^{r-2j-2}_{\Delta+j_1+1} .
\tag{2.26}
\]
链条两端会发生什么？显然，链条不能任意向右延伸，因为共形维数在那个方向上递减。另一方面，右端的 $B_1$ 型多重态具有零态 $[\frac12,j_2]^{r-2j-3}_{\Delta+j+\frac12}$。这个零态有 $j_1\neq0$，并且容易验证它位于幺正性界之下，因此根本谈不上以它作为原初态构造多重态。这样，重组链在两个方向上都终止。

% [原文 p.17]
为着重组的用途，把末端的 $B_1$ 型多重态看作 $j_1=-\frac12$ 的 $A_1$ 型多重态是方便的，即
\[
[0,j_2]^{r}_{\Delta=-\frac32 r}\equiv[-\tfrac12,j_2]^{r+1}_{\Delta=-\frac32 r-\frac12} .
\tag{2.27}
\]
超共形指标需要捕捉模掉重组之后的全部谱信息；由于上述链条中任意两个相邻多重态都能重组，指标应当计入
\[
N_{r'=r-2j_1,j_2}=\sum_{p=-1}^{k}(-1)^p\,n\big([\tfrac{p}{2},j_2]^{r'+p}_{\Delta=2+p-\frac32(r'+p)}\big) .
\tag{2.28}
\]
确实，$A_1$ 型多重态 $[j_1,j_2]^r_\Delta$ 的指标正是 [25]，
\[
\sI_{[j_1,j_2]^r_\Delta}=(-)^{2(j_1+j_2)+1}\frac{(pq)^{\frac12(-r+2j_1+2)}\chi_{j_2}\big(\sqrt{q/p}\big)}{(1-p)(1-q)} .
\tag{2.29}
\]
若我们在重组链中的一族多重态上求值，它恰好计入 $N_{r'=r-2j_1,j_2}$。这是因为，除符号交替之外，链条中所有多重态的指标都相同。

$\sN=1$ 超共形代数为 $\SU(2,2|1)$，该代数中单个超荷的对易子为 $\SU(2,1)$。指标必须是该代数的特征标。确实，该代数的两个 Cartan 生成元在 $\SU(2,2|1)$ 中的嵌入为 $r-2j_1$ 与 $j_2$。这与重组规则的预期一致。

\subsection{中场休息：超对称拉氏量概览}

\paragraph{$\sN=0$ 拉氏量：} 一条经验法则是：超对称越多，指定拉氏量所需的数据越少。例如，若想构造标量、费米子与规范场的普遍拉氏量。首先必须指定物质场在规范对称性下的荷（或更一般地，表示）。除规范相互作用外，物质还可以带有自身的相关相互作用，例如 Yukawa 项与四次标量项，并以保持规范对称性的方式出现。最后，费米子的数目及其表示需要仔细选取，以使得规范对称性无反常。

\paragraph{$\sN=1$ 拉氏量：} 有了 $\sN=1$ 超对称，分类多少简化了一些（尽管远不及 $\sN=2$ 那样简单）。标量与费米子在一个多重态中，即手征多重态；规范场与规范微子在一个多重态中，即矢量多重态。指定手征多重态在规范对称性下荷的表示，就完全固定了手征—矢量相互作用。对 $\sN=0$，物质相互作用有两类且彼此独立；但有了 $\sN=1$ 超对称，它们彼此关联，并可写成一个超势 $W(\Phi_i)$。同样，需要确保物质内容使得规范对称性无反常。

% [原文 p.18]
\paragraph{$\sN=2$ 拉氏量：} $\sN=2$ 拉氏量用 $\sN=2$ 超多重态（由一对手征多重态组成）与 $\sN=2$ 矢量多重态（由 $\sN=1$ 矢量多重态与 $\sN=1$ 手征多重态组成，二者都在规范群的伴随表示下变换）来构造。具有 $\sN=2$ 超对称的拉氏量比 $\sN=1$ 的容易指定得多。只需指定矢量—物质的耦合；物质—物质之间的耦合被 $\sN=2$ 超对称禁止。前者只需通过指定超多重态在规范群下的表示来给出。这使得可以用称为箭图的便利图解记号来描述 $\sN=2$ 理论。此外，由于物质成对地处于共轭表示中，$\sN=2$ 理论的规范对称性总是无反常的。因此以这种方式指定的拉氏量总是好的拉氏量（不过它可能不是渐近自由的）。当规范耦合的单圈 beta 函数为零时，规范耦合充当恰当边缘形变。在那种情形下，相互作用超共形理论由复化规范耦合参数化。

\paragraph{$\sN=3$ 拉氏量：} 不存在不具有 $\sN=4$ 对称性的 $\sN=3$ 拉氏量。不过近年来发现了一些非拉氏量的 $\sN=3$（但非 $\sN=4$）理论 [26–28]。$\sN=1$ 在 $\sN=3$ 中的对易子秩为二。因此可以再开启两个逸度，使指标成为四个变量的函数。关于 $\sN=3$ 理论指标的讨论见 [29]。

\paragraph{$\sN=4$ 拉氏量：} $\sN=4$ 拉氏量由伴随表示中的单个矢量多重态构造。它总是共形的，复化规范耦合参数化超共形理论的空间。$\sN=4$ 矢量多重态由伴随表示中的 $\sN=2$ 矢量多重态与 $\sN=2$ 超多重态组成。用 $\sN=1$ 多重态来说，它由一个 $\sN=1$ 矢量多重态与三个在伴随表示中变换的 $\sN=1$ 手征多重态组成。$\sN=1$ 在 $\sN=4$ 中的对易子秩为二\footnote{它并非人们天真猜测的三，原因在于 R 对称性是 $\SU(4)$ 而非 $\SU(4)\times\U(1)$。对 $\sN=1,2,3$，R 对称性是 $\SU(N)\times\U(1)$。}。因此可以再开启两个逸度，使指标成为四个变量的函数。

如果超共形代数的秩为 $r$，则最一般的指标是 $r-2$ 个逸度的函数。这是因为指标只计入 $\delta=0$ 的态，而我们只能对与 $Q$ 对易的荷开启逸度。
% [原文 p.19]
这两个条件把变量数减少 2。

% [Tutor Note: 术语：表 2 中的「字母」是单字母，$\sI$ 列给出该字母对单字母指标分子（或分母塔）的贡献，而不是完整指标。]
% [Tutor Note: 易错：表 2 只列 $\delta=0$ 的字母，这与表 1 不同；$\delta\neq0$ 的字母对指标无贡献，故不列出。]
% [Tutor Note: 易错：$\chi_{\adj}$ 中的加性 $N$ 来自 Cartan 生成元（零权），不要漏掉，否则 $i_V$ 与 (2.17) 不符。]
% [Tutor Note: 动机：指标在规范耦合变化下不变，因此可在弱耦合（$g=0$）下计算，再解析延拓到强耦合——这是指标能给出非平凡检验的根本原因。]
% [Tutor Note: 易错：(2.18) 是被积函数在规范群上的投影；$1/|W|$ 与 $\Delta(a_i)$ 一起实现 Weyl 投影，$\Delta$ 的幂次是 $+1$。]
% [Tutor Note: 易错：(2.20) 原文照录，但测度应为 $\prod_i \dd a_i/(2\pi\ii)$，且单独的 $\Gamma(pqa^\alpha)$ 应为 $\prod_\alpha\Gamma(pqa^\alpha)$；已登记 ERRATA。]
% [Tutor Note: 术语：Spiridonov beta 函数恒等式 (2.22) 即数学中的椭圆 beta 积分，$\SU(2)$ 三味指标的闭式正是它。]
% [Tutor Note: 动机：(2.21) 中 $z^{\pm1}$ 是简写 $\{z,z^{-1}\}$，$z^{\pm2}=\{z^2,z^{-2}\}$，不要漏掉任何一支。]
% [Tutor Note: 易错：Seiberg 对偶两边指标相等，这是对偶的非平凡检验；但它不证明对偶，只检验 BPS 谱。]
% [Tutor Note: 术语：长/短多重态、重组、零态、幺正性界——短多重态在重组下可与长多重态互换，指标对重组不敏感。]
% [Tutor Note: 易错：(2.24) 原文正文说短多重态是 $\{|\chi\rangle\}$ 而方程写 $Q|\phi\rangle=0$，符号不一致；已登记 ERRATA。]
% [Tutor Note: 易错：(2.26) 末项指数原文作 $r-2j-2$、下段作 $r-2j-3$，其中 $j$ 应为 $j_1$；译文照录，已登记 ERRATA。]
% [Tutor Note: 动机：(2.28) 的 $(-1)^p$ 正是长短多重态交替的符号，求和后把链条的贡献约化为一个数。]
% [Tutor Note: 术语：$\sN=1$ 指标是 $\SU(2,1)$ 的特征标，两个 Cartan 为 $r-2j_1$、$j_2$；这解释了 (2.29) 的形式。]
% [Tutor Note: 文献：[25] 给出 $A_1$ 型多重态指标的显式形式；[23] 最早讨论重组规则；[12] 为本文采用的叙述。]
% [Tutor Note: 动机：$\sN=2,3,4$ 的指标分别依赖 3、4、4 个变量，因为对易子的秩随超对称增加而增加。]


\section{$\sN=2$ 理论}
% sec3a.tex — c06｜原文 p19–22（p19 从 §3 标题后起）｜§3 N=2 理论、§3.1 超对称极限、§3.2 S 对偶、§3.3 开头
% 编号公式 (3.1)–(3.19) 共 19 条；编号环境：习题 3.1、3.2。
% 注意：p19 顶端两句（「这两个条件把变量数减少 2」）已在 sec2b 末尾，本块不重复。
% [原文 p.19]
$\sN=2$ 超共形代数由两个超荷
$\{\mathcal{Q}_{I\alpha},\mathcal{S}^{I\alpha}\equiv\mathcal{Q}^{\dagger I\alpha},\tilde{\mathcal{Q}}^I_{\dot\alpha},\tilde{\mathcal{S}}^{\dot\alpha}_I\equiv\tilde{\mathcal{Q}}^{\dagger\dot\alpha}_I\}$
组成，其中 $I=1,2$。除共形代数外，玻色子代数还包括 R 对称性 $\SU(2)_R\times\U(1)_r$。超荷 $\mathcal{Q}_{I\alpha}$ 在 $\SU(2)$ 下按二重态变换，该 $\SU(2)$ 由指标 $I$ 显式标出。在 $\U(1)_r$ 下，$\mathcal{Q}_{I\alpha}$ 的荷为 $-1$，而 $\tilde{\mathcal{Q}}^I_{\dot\alpha}$ 的荷为 $1$。指标 $\alpha$ 与 $\dot\alpha$ 分别是 $\SU(2)_1$ 与 $\SU(2)_2$ 的指标，而 $\SU(2)_1\times\SU(2)_2=\mathrm{Spin}(4)$ 是共形群的转动子群。所有 $\mathcal{Q}$ 与 $\tilde{\mathcal{Q}}$ 的共形维数都是 $1/2$，所有 $\mathcal{S}$ 与 $\tilde{\mathcal{S}}$ 的共形维数都是 $-1/2$。超共形代数中最重要的（反）对易关系为
\[
\{\mathcal{Q}_{I\alpha},\mathcal{Q}^{\dagger J\beta}\}
=\Delta\delta_\alpha^\beta\delta_I^J+2M_\alpha^\beta\delta_I^J-2R_I^J\delta_\alpha^\beta+\frac12 r\,\delta_\alpha^\beta\delta_I^J ,
\tag{3.1}
\]
\[
\{\tilde{\mathcal{Q}}^I_{\dot\alpha},\tilde{\mathcal{Q}}^{\dagger\dot\beta}_J\}
=\Delta\delta_{\dot\alpha}^{\dot\beta}\delta_J^I+2\tilde M_{\dot\alpha}^{\dot\beta}\delta_J^I-2R_J^I\delta_{\dot\alpha}^{\dot\beta}-\frac12 r\,\delta_{\dot\alpha}^{\dot\beta}\delta_J^I .
\tag{3.2}
\]
不失一般性，我们取 $\mathcal{Q}_{1-}$ 来计算指标。显式地写出下式是有用的：
\[
\{\mathcal{Q}_{1-},(\mathcal{Q}_{1-})^\dagger\}=\Delta-2j_1-2R+\frac12 r .
\tag{3.3}
\]
选定了超荷 $\mathcal{Q}_{1-}$，就确定了一个 $\sN=1$ 子代数。它由
$\{\mathcal{Q}_{1\alpha},\mathcal{S}^{1\alpha}\equiv\mathcal{Q}^{\dagger1\alpha},\tilde{\mathcal{Q}}^1_{\dot\alpha},\tilde{\mathcal{S}}^{\dot\alpha}_1\equiv\tilde{\mathcal{Q}}^{\dagger\dot\alpha}_1\}$
生成。与式 (2.2) 比较，
\[
r_{\sN=1}=\frac23\Big(-2R+\frac12 r\Big) .
\tag{3.4}
\]
此外，所有超荷都与 $R+\frac12 r$ 对易。由于只有超荷与 R 对称性有非平凡的对易关系，$R+\frac12 r$ 就是相对于整个 $\sN=1$ 代数的对易子。我们通过开启关于该荷的逸度来定义指标，
\[
\sI(p,q,t)\equiv\Tr(-1)^F p^{\frac13(\Delta+j_1)+j_2}q^{\frac13(\Delta+j_1)-j_2}x^{(R+\frac12 r)},\qquad
\delta=\Delta-2j_1-2R+\frac12 r ,
\tag{3.5}
\]
\paragraph{超多重态的指标} $\sN=2$ 超多重态由一对手征多重态组成。一个手征多重态与另一个手征多重态的复共轭一起，在 $\SU(2)_R$ 对称性下按二重态 $q_I$ 变换，并且在整体对称性下具有相同的表示。这意味着超多重态中的两个手征多重态按整体对称性的共轭表示变换。我们来看它们对指标的贡献。

% [原文 p.20]
手征多重态 $q_1$ 的 $r_{\sN=1}=\frac23(1+0)$，另一个手征多重态 $\bar q_2$ 也是如此。此外，它们在 $R+\frac12 r=-\frac12$ 这个对易子下的荷相同。因此超多重态的单字母指标为
\[
\begin{aligned}
i_H(a_i;p,q,x)&=\frac{(pq)^{\frac13}x^{\frac12}\chi_R(a_i)-(pq)^{\frac23}x^{-\frac12}\chi_{\bar R}(a_i)}{(1-p)(1-q)}
+\frac{(pq)^{\frac13}x^{\frac12}\chi_{\bar R}(a_i)-(pq)^{\frac23}x^{-\frac12}\chi_R(a_i)}{(1-p)(1-q)}\\
&=\frac{(pq)^{\frac13}x^{\frac12}-(pq)^{\frac23}x^{-\frac12}}{(1-p)(1-q)}\Big(\chi_R(a_i)+\chi_{\bar R}(a_i)\Big) .
\end{aligned}
\tag{3.6}
\]
单个手征多重态 $q_1$ 有时称为半超多重态。超多重态的贡献可以理解为半超多重态取遍表示 $R$ 与 $\bar R$ 时的贡献。

\paragraph{矢量多重态的指标} $\sN=2$ 矢量多重态由一个 $\sN=1$ 手征多重态 $\phi$ 与一个 $\sN=1$ 矢量多重态 $V$ 组成。该手征多重态在 $\SU(2)_R$ 下中性，但在 $\U(1)_r$ 下的荷为 2。有效的 $r_{\sN=1}=\frac23(0+\frac12)=\frac23$，在对易子下的荷为 $R+\frac12 r=1$。当然，$\sN=1$ 矢量多重态不能带有任何整体对称性的荷。因此，$\sN=2$ 矢量多重态的指标为
\[
i_V(a_i;p,q,x)=\Big(\frac{(pq)^{\frac13}x^{-1}-(pq)^{\frac23}x}{(1-p)(1-q)}+\frac{-p-q+2pq}{(1-p)(1-q)}\Big)\chi_{\adj}(a_i) .
\tag{3.7}
\]
\subsection{超对称极限}
可以取超共形指标的极限，使得并非全部短多重态、而只有某些特殊的短多重态对指标有贡献。$\sN=2$ 指标的这类超对称极限已在 [30] 中分类并命名。为便于分析，采用变量 $t$，作替换 $x\to(pq)^{-\frac23}t$。用这些变量，指标为
\[
\sI(p,q,t)=\Tr(-1)^F p^{\frac13(\Delta+j_1)+j_2-\frac23(R+\frac12 r)}q^{\frac13(\Delta+j_1)-j_2-\frac23(R+\frac12 r)}t^{R+\frac12 r} .
\tag{3.8}
\]
利用 $\delta=\Delta-2j_1-2R+\frac12 r=0$，它可以改写为
\[
\begin{aligned}
\sI(p,q,t)&=\Tr(-1)^F p^{\frac12(\Delta+2j_2-2R-\frac12 r)}q^{\frac12(\Delta-2j_2-2R-\frac12 r)}t^{R+\frac12 r},\\
&=\Tr(-1)^F p^{\frac12\{\tilde{\mathcal{Q}}_{1\dot+},(\tilde{\mathcal{Q}}_{1\dot+})^\dagger\}}q^{\frac12\{\tilde{\mathcal{Q}}_{1\dot-},(\tilde{\mathcal{Q}}_{1\dot-})^\dagger\}}t^{R+\frac12 r} .
\end{aligned}
\tag{3.9}
\]
于是半超多重态与矢量多重态的指标变为
\[
i_{\frac12 H}(p,q,t)=\frac{\sqrt t-pq/\sqrt t}{(1-p)(1-q)},\qquad
i_V(p,q,t)=\frac{pq/t-t}{(1-p)(1-q)}+\frac{-p-q+2pq}{(1-p)(1-q)} .
\tag{3.10}
\]
把指标写成式 (3.9) 的形式，其好处是出现的荷显为非负。这允许把 $p$ 或 $q$ 或二者同时取为零。如下文所见，这些极限计入具有增强超对称性的算符。
% [原文 p.21]
\begin{exercise}[3.1]
定义 $\delta_\alpha=\{Q_\alpha,(Q_\alpha)^\dagger\}$，并相应地定义 $\tilde\delta_{\dot\alpha}$，用荷 $\delta_\alpha,\tilde\delta_{\dot\alpha}$ 构造指标。这样选择荷的好处是它们是正定的，因而相应的逸度可以取为零。
\end{exercise}
\paragraph{Macdonald 指标：} 对应 $p\to0$ 的极限。在此极限下，$\{\tilde{\mathcal{Q}}_{1\dot+},(\tilde{\mathcal{Q}}_{1\dot+})^\dagger\}=0$ 的态有贡献。这意味着它们除了被 $\mathcal{Q}_{1-}$$\,$熄灭外，还被 $\tilde{\mathcal{Q}}_{1\dot+}$ 熄灭。在此极限下单字母指标变为
\[
i_{\frac12 H}(q,t)=\frac{\sqrt t}{1-q},\qquad i_V(q,t)=\frac{-t-q}{1-q} .
\tag{3.11}
\]
\paragraph{Hall–Littlewood 指标：} 这是指标的 $p\to0,\ q\to0$ 极限。在此极限下，算符除了被 $\mathcal{Q}_{1-}$ 熄灭外，还被 $\tilde{\mathcal{Q}}_{1\dot+}$ 与 $\tilde{\mathcal{Q}}_{1\dot-}$ 熄灭。条件 $\delta=0$、$\Delta+2j_2-2R-\frac12 r=0$、$\Delta-2j_2-2R-\frac12 r=0$ 可以化简为
\[
j_2=0,\qquad j_1=r,\qquad \Delta=2R+r .
\tag{3.12}
\]
此极限下的单字母指标为
\[
i_{\frac12 H}(t)=\sqrt t,\qquad i_V(t)=-t .
\tag{3.13}
\]
\paragraph{Schur 指标：} Schur 极限是指标的 $q=t$ 极限。
\[
\sI(p,q)=\Tr(-1)^F p^{\frac12(\Delta+2j_2-2R-\frac12 r)}q^{\frac12(\Delta-2j_2-2R-\frac12 r)+(R+\frac12 r)}
\tag{3.14}
\]
指标中出现的荷是 $\Delta+2j_2-2R-\frac12 r$ 与 $\Delta-2j_2+\frac12 r$。由构造，这些荷与 $\mathcal{Q}_{1-}$ 对易，因为我们计算的正是关于 $\mathcal{Q}_{1-}$ 的指标。有趣的是，这些荷也与 $\tilde{\mathcal{Q}}_{1\dot+}$ 对易。因此，指标只从 $\{\tilde{\mathcal{Q}}_{1\dot+},(\tilde{\mathcal{Q}}_{1\dot+})^\dagger\}=0$ 的态获得贡献。于是指标自动与 $p$ 无关，迹公式化简为
\[
\sI(q)=\Tr(-1)^F q^{\Delta-R} .
\tag{3.15}
\]
此极限下的单字母指标为
\[
i_{\frac12 H}(q)=\frac{\sqrt q}{1-q},\qquad i_V(q)=\frac{-2q}{1-q} .
\tag{3.16}
\]
超对称极限使人得以探测 $\sN=2$ 超共形代数的额外超对称扇区。
% [原文 p.22]
\subsection{$\sN=2$ 超 QCD 的 S 对偶}
一大类 $\sN=2$ 规范理论（称为 class S）的非平凡对偶由 Gaiotto 发现 [31]。这类对偶中最简单的一个——已由 Seiberg 与 Witten 的工作 [32] 所知——是带八个半超多重态的 $\SU(2)$ 规范理论。这些半超多重态按 $\SO(8)$ 的矢量表示变换。我们只看子群 $\SO(8)\supset\SO(4)\times\SO(4)=\SU(2)_a\times\SU(2)_b\times\SU(2)_c\times\SU(2)_d$。$\SO(8)$ 按 $\SU(2)^4$ 的表示分解为 $\mathbf{8}_v=(2_a\otimes2_b)\oplus(2_c\otimes2_d)$。该表示的特征标为
\[
\Big(a+\frac1a\Big)\Big(b+\frac1b\Big)+\Big(c+\frac1c\Big)\Big(d+\frac1d\Big) .
\tag{3.17}
\]
有了这些信息，我们立即可以写下指标的积分公式。该理论的指标为
\[
\sI=\kappa\,\Gamma\Big(\frac{pq}{t}\Big)\oint\frac{\dd z}{4\pi\ii z}\frac{\Gamma\big(\frac{pq}{t}z^{\pm2}\big)}{\Gamma(z^{\pm2})}\Gamma\big(\sqrt t\,z^{\pm1}a^{\pm1}b^{\pm1}\big)\Gamma\big(\sqrt t\,z^{\pm1}c^{\pm1}d^{\pm1}\big) .
\tag{3.18}
\]
这里我们用了简写 $\Gamma(z^\pm)=\Gamma(z)\Gamma(z^{-1})$ 等等。

已知该理论具有三重性（triality）。若把半超多重态的矢量表示换成旋量表示或共轭旋量表示，则得到同一理论，但耦合不同。由于指标对耦合不敏感，我们预期指标在这些表示改变下不变。旋量与共轭旋量按 $\SU(2)^4$ 分解为
\[
8_s=(2_a\otimes2_c)\oplus(2_b\otimes2_d),\qquad
8_c=(2_a\otimes2_d)\oplus(2_b\otimes2_c) .
\tag{3.19}
\]
这意味着指标必须在交换 $b\leftrightarrow c$ 与 $b\leftrightarrow d$ 下不变。除这些非平凡的不变性外，指标显然在平凡交换 $a\leftrightarrow b$ 与 $c\leftrightarrow d$ 下不变。

在数学文献中，这些不变性被发现 [33] 的时间与物理学中预期它们 [34] 的时间大致相同。作为物理学家，我们可以通过令 $p=xy$、$q=x/y$ 并按 $x$ 的幂展开来快速检验这些不变性。
\begin{exercise}[3.2]
完成此事。
\end{exercise}
\subsection{class S 对偶与 TQFT}
根据 Gaiotto 的工作，已知该理论是通过把六维秩 1 的 $(2,0)$ 理论紧化在带四个穿刺点的球面上得到的。参数化该曲面的复结构映射为复化规范耦合。详细讨论见 Bruno LeFloch 的讲义。此外，$\SU(2)^4$ 整体对称性中的每个 $\SU(2)$ 因子对应一个穿刺点。

% [Tutor Note: 术语：本节超荷用花体 $\mathcal{Q}$（原文如此），而习题 3.1 与 §2 用正体 $Q$；不要因字形不同而改符号。]
% [Tutor Note: 易错：(3.1) 与 (3.2) 中 $R$ 的指标次序不同（$R_I^J$ 与 $R_J^I$），且 $r$ 项符号相反（$+\frac12r$ 与 $-\frac12r$），这是 $\U(1)_r$ 手征性的体现。]
% [Tutor Note: 易错：本节 $\delta$ 含 $-2R$ 项，与 §2 的 $\delta=\Delta-2j_1+\frac32r$ 不同；换了超荷，$\delta$ 随之改变。]
% [Tutor Note: 动机：(3.4) 把 $\sN=2$ 的 R 荷翻译成 $\sN=1$ 的 $r$ 荷，这样才能对每个 $\sN=1$ 多重态套用 §2 的单字母公式。]
% [Tutor Note: 术语：hypermultiplet 译「超多重态」，half-hypermultiplet 译「半超多重态」；二者都不是 $\sN=1$ 的独立多重态。]
% [Tutor Note: 易错：(3.6) 第一行两个分式中 $\chi_R$ 与 $\chi_{\bar R}$ 的次序相反，第二个分式来自 $\bar q_2$；不要抄成同一个分式。]
% [Tutor Note: 动机：(3.9) 把指数写成超荷反对易子的形式，使荷显为非负，这是取 $p,q\to0$ 极限（Macdonald/Hall–Littlewood）的前提。]
% [Tutor Note: 易错：(3.9) 中 $p,q$ 的指数分别是两个反对易子的一半，不是 (3.5) 的 $\frac13(\Delta+j_1)\pm j_2$；两者只在 $\delta=0$ 时等价。]
% [Tutor Note: 物理：三个极限分别对应 $p\to0$（Macdonald）、$p,q\to0$（Hall–Littlewood）、$q=t$（Schur），各自只保留增强超对称的扇区。]
% [Tutor Note: 易错：Schur 极限是令 $q=t$，不是令 $t=0$；结果 (3.15) 只依赖一个变量 $q$。]
% [Tutor Note: 物理：(3.12) 的三个条件 $j_2=0$、$j_1=r$、$\Delta=2R+r$ 是 Hall–Littlewood 极限下存活态的 BPS 条件。]
% [Tutor Note: 术语：triality 译「三重性」；$\mathbf{8}_v,\mathbf{8}_s,\mathbf{8}_c$ 的分解正是 $\SO(8)$ 三重性的群论体现。]
% [Tutor Note: 易错：(3.18) 中被积函数含两个 $\Gamma$ 因子（分别含 $a,b$ 与 $c,d$），不是四个；$z^{\pm2}$ 是 $z^2,z^{-2}$ 的简写。]
% [Tutor Note: 文献：[31] Gaiotto 的 class S 构造；[32] Seiberg–Witten；[33][34] 数学与物理几乎同时发现这些不变性。]
% [Tutor Note: 物理：指标对耦合不敏感，因此 S 对偶两边指标必须相等；这是检验对偶的强约束，但不是对偶的证明。]
% [Tutor Note: 语病：原文 p.22 「the a puncture」冠词衍余，应为「a puncture」；照译并登记 ERRATA。]

% sec3b.tex — c07｜原文 p23–25｜§3.3 class S 对偶与 TQFT（续）、Schur 极限、推广到更高秩
% 编号公式 (3.20)–(3.35) 共 16 条；无编号环境、无表格。
% [原文 p.23]
在退化极限下，四穿刺球面退化并分裂成三个三穿刺球面。根据 Gaiotto，一个三穿刺球面是相对于三个穿刺点处 $\SU(2)$ 的三重基础表示中的半超多重态。带三个穿刺点的球面没有任何模，这与自由半超多重态理论没有任何耦合相符。可以取三种不同的退化极限，即分别把 $a$ 与 $b$、$a$ 与 $c$、或 $a$ 与 $d$ 靠近。在每个退化极限下我们都得到一个弱耦合理论，但其中半超多重态分别处于 $\SO(8)$ 整体对称性的 $\mathbf{8}_v$、$\mathbf{8}_s$ 与 $\mathbf{8}_c$ 表示。只看 $\SU(2)^4$ 子群，该对偶可以看作一种交叉对称性。

更一般地，由于指标与规范耦合无关，因而与 Riemann 曲面的复结构无关，它必须由某个拓扑场论来计算。抽象地说，一个拓扑场论可以通过给出三点函数与两点函数（即传播子）来指定。

让我们回顾这是如何运作的。我们把三穿刺球面的指标参数化为 $\sI(a,b,c)$，其中 $a,b,c$ 是三个 $\SU(2)$ 对称性 Cartan 的逸度。另一方面，我们可以容易地写下与圆柱相对应的「传播子」，
\[
\eta(a,b)=\Delta(a)\sI_V(a)\,\delta(a,b^{-1}) ,
\tag{3.20}
\]
其中 $\Delta(a)$ 是 Haar 测度，$\sI_V(a)$ 是矢量多重态的指标，后者已知显式形式。一般 class S 理论的指标可以用这些基本组元的指标表示。最简单的例子是：用一个圆柱把两个三穿刺球面粘起来，就得到四穿刺球面的指标，
\[
\sI(a,b,c,d)=\oint\frac{\dd z}{2\pi\ii}\oint\frac{\dd x}{2\pi\ii}\sI_H(a,b,z)\,\eta(z,x)\,\sI_H(z,c,d)
\tag{3.21}
\]
\[
=\oint\dd z\,\Delta(z)\,\sI_H(a,b,z)\,\sI_V(z)\,\sI_H(z,c,d) .
\tag{3.22}
\]
若把指标按方便的函数基 $\{f^\alpha(a)\}$ 展开，该基以 $\SU(2)$ 表示 $\{\alpha\}$ 标号\footnote{对 type A 理论，$\{f^\alpha(a)\}$ 是其自变量的对称函数，这些自变量是与 $\SU(k)$ 的 Cartan 生成元对偶的逸度。更一般地，对 type D 与 E 理论，$\{f^\alpha(a)\}$ 在相应的 Weyl 群下不变。}，我们就可以给每个三穿刺球面指定「结构常数」$C_{\alpha\beta\gamma}$，给每个传播子指定度规 $\eta^{\alpha\beta}$，
\[
\sI(a,b,c)=\sum_{\alpha,\beta,\gamma}C_{\alpha\beta\gamma}\,f^\alpha(a)\,f^\beta(b)\,f^\gamma(c)
\tag{3.23}
\]
\[
\eta^{\alpha\beta}=\oint\frac{\dd a}{2\pi\ii}\oint\frac{\dd b}{2\pi\ii}\eta(a,b)\,f^\alpha(a)\,f^\beta(b) .
\tag{3.24}
\]
指标在曲面的不同分解方式下不变，等价于
% [原文 p.24]
说 $C_{\alpha\beta\gamma}$ 与 $\eta^{\alpha\beta}$ 定义了一个二维拓扑场论\footnote{我们在这个词的用法上有些不严格。按 Atiyah 的公理化定义，TQFT 应具有有限维状态空间，而在我们的情形状态空间是无穷维的。具有无穷维状态空间的二维拓扑理论中理解得最清楚的例子是二维 Yang–Mills 理论的零面积极限 [35,36]（综述见 [37]）。幸运的是，与指标相伴的二维拓扑理论与二维 Yang–Mills 密切相关。}。关键性质是结合性，
\[
C_{\alpha\beta\zeta}\,C^{\zeta}_{\ \gamma\delta}=C_{\alpha\gamma\zeta}\,C^{\zeta}_{\ \beta\delta} ,
\tag{3.25}
\]
其中指标用度规 $\eta^{\alpha\beta}$ 上升、用逆度规 $\eta_{\alpha\beta}$ 下降。

很自然的选择是让完备函数集 $\{f^\alpha(a)\}$ 在传播子中出现的测度下正交归一，
\[
\oint\frac{\dd a}{2\pi\ii}\Delta(a)\sI_V(a)\,f^\alpha(a)f^\beta(a)=\delta^{\alpha\beta} .
\tag{3.26}
\]
于是度规 $\eta^{\alpha\beta}$ 是平庸的，
\[
\eta^{\alpha\beta}=\delta^{\alpha\beta} .
\tag{3.27}
\]
条件 (3.26) 仍留下相当大的自由度，因为无穷多个彼此由正交变换联系起来的函数基都满足它。真正的简化出现在我们能找到一个显式基 $\{f^\alpha(a)\}$，使得结构常数是对角的，
\[
C_{\alpha\beta\gamma}\neq0\quad\to\quad\alpha=\beta=\gamma .
\tag{3.28}
\]
此时结合性 (3.25) 自动成立。对满足 (3.25) 的结构常数，总可以找到使它们对角的基。困难在于以具体形式描述该基。

一般而言，传播子中出现的测度很复杂，没有现成的显式正交归一函数集。我们觉得考虑如下拟设（ansatz）非常有用
\[
f^\alpha(a)=\mathcal{K}(a)\,P^\alpha(a) ,
\tag{3.29}
\]
其中 $\mathcal{K}(a)$ 是某个函数。显然，由 (3.26)，函数 $\{P^\alpha(a)\}$ 在新测度 $\hat\Delta(a)$ 下正交归一，
\[
\oint\frac{\dd a}{2\pi\ii}\hat\Delta(a)\,P^\alpha(a)P^\beta(a)=\delta^{\alpha\beta},\qquad
\hat\Delta(a)\equiv\sI_V(a)\mathcal{K}(a)^2\Delta(a) .
\tag{3.30}
\]
回忆 $\Delta(a)$ 是 Haar 测度。问题的关键在于巧妙地选取 $\mathcal{K}(a)$，使得 $\hat\Delta(a)$ 是简单的已知测度，且正交归一基 $\{P^\alpha(a)\}$ 是满足 (3.28) 的显式函数集。

一旦已知对角基 $\{f^\alpha(a)\}$ 与结构常数 $C_{\alpha\alpha\alpha}$，就容易计算与亏格 $g$、带 $s$ 个穿刺点的曲面相伴的 SCFT 的指标。这样的曲面可以通过粘合 $2g-2+s$ 个三穿刺球面得到，于是
\[
\sI_{g,s}(a_1,a_2,\ldots,a_s)=\sum_\alpha\big(C_{\alpha\alpha\alpha}\big)^{2g-2+s}\prod_{I=1}^{s}f^\alpha(a_I) .
\tag{3.31}
\]
\paragraph{Schur 极限} 在 Schur 极限下最容易辨认正交归一函数基。在此极限下，
\[
\sI_V(a)=PE\Big[\frac{-2q}{1-q}(a^2+1+a^{-2})\Big],\qquad
\sI_H(a,b,c)=PE\Big[\frac{\sqrt q}{1-q}\Big(a+\frac1a\Big)\Big(b+\frac1b\Big)\Big(c+\frac1c\Big)\Big] .
\tag{3.32}
\]
% [原文 p.25]
若取 $\mathcal{K}(a)=1/\sqrt{\sI_V(a)}$，则 $\hat\Delta(a)=\Delta(a)$。关于标准 Haar 测度正交归一的函数基是特征标，也称为 Schur 多项式。确实，三点函数在该基下是对角的 [38]，即
\[
\sI_H(a,b,c)=\frac{C(q)}{\sqrt{\sI_V(a)\sI_V(b)\sI_V(c)}}\sum_\alpha\frac{1}{\dim_q\alpha}\chi_\alpha(a)\chi_\alpha(b)\chi_\alpha(c) .
\tag{3.33}
\]
由于与该指标极限相伴的正交多项式是 Schur 多项式，该极限称为 Schur 极限。

在 Hall–Littlewood 与 Macdonald 极限中，出现的正交多项式分别是 Hall–Littlewood 与 Macdonald 多项式。对完整指标，$P_\alpha(a)$ 不是多项式，而是与 Ruijsenaars–Schnider 可积模型相关的复杂函数。该模型位于 Calogero–Moser 层级的顶端。

\paragraph{推广到更高秩} 如前所述，$\sN=2$ 的 $\SU(2)$ 规范理论由六维秩 1 的 $(2,0)$ SCFT 紧化得到。若把更高秩的 $(2,0)$ 理论紧化在 Riemann 曲面上，会得到类似的网络，但与秩 1 情形不同，这些理论大多没有拉氏量描述。特别地，对秩 $N$，与三穿刺球面相伴的理论称为 $T_N$ 理论。$T_2$ 理论只是自由超多重态的理论，而 $T_{N>2}$ 是非平庸的强耦合 SCFT。这使人无法走计算指标的直接路线。然而，由于拓扑三点函数与两点函数的表达式具有可自然提升到更高秩的形式，人们就得到了整个强耦合理论族的超共形指标的自然猜想。同样，在 Schur 极限下举例最简单。由于 $T_2$ 理论的指标为
\[
\sI_{T_2}(a,b,c)=\frac{C_2(q)}{\sqrt{\sI_V(a)\sI_V(b)\sI_V(c)}}\sum_\alpha\frac{1}{\dim_q\alpha}\chi_\alpha(a)\chi_\alpha(b)\chi_\alpha(c) ,
\tag{3.34}
\]
这一形式到一般 $N$ 的提升将是
\[
\sI_{T_N}(a_i,b_i,c_i)=\frac{C_N(q)}{\sqrt{\sI_V(a_i)\sI_V(b_i)\sI_V(c_i)}}\sum_R\frac{1}{\dim_q R}\chi_R(a_i)\chi_R(b_i)\chi_R(c_i) ,
\tag{3.35}
\]
其中 $R$ 遍历 $\SU(N)$ 的所有表示，$a_i,b_i,c_i$ 是与整体对称性 $\SU(N)_i$ 相伴的逸度。

对 Macdonald 极限 [38] 乃至完整指标 [39] 也已作出类似的猜想。

% [Tutor Note: 物理：class S 理论即六维 $(2,0)$ 理论紧化在 Riemann 曲面上；三穿刺球面是构建一切的基本积木。]
% [Tutor Note: 易错：(3.20) 传播子中是 $\delta(a,b^{-1})$ 而不是 $\delta(a,b)$；$\SU(2)$ 特征标以 $a\leftrightarrow a^{-1}$ 配对。]
% [Tutor Note: 动机：指标与耦合无关，故只依赖曲面的复结构，因此必是拓扑场论（TQFT）的量——这是本节的中心思想。]
% [Tutor Note: 易错：(3.21) 与 (3.22) 各占一个编号：前者保留 $\eta$，后者把 $\eta(z,x)$ 代入并使用 $\delta(z,x^{-1})$ 消去一个积分。]
% [Tutor Note: 术语：结构常数 $C_{\alpha\beta\gamma}$、度规 $\eta^{\alpha\beta}$、结合性 (3.25)，正是 Atiyah 二维 TQFT 公理的具体化。]
% [Tutor Note: 易错：(3.26) 的正交归一测度是 $\Delta(a)\sI_V(a)$，不是单独的 Haar 测度；后者只在 Schur 极限的 $\hat\Delta=\Delta$ 特例中出现。]
% [Tutor Note: 易错：(3.24) 定义 $\eta^{\alpha\beta}$（指标在上），(3.25) 用 $\eta^{\alpha\beta}$ 上升、$\eta_{\alpha\beta}$ 下降；原文把二者称作 metric 与 inverse metric。]
% [Tutor Note: 物理：对角化结构常数 (3.28) 使结合性自动满足，问题化为「找基」；Schur 极限下这组基就是 Schur 多项式。]
% [Tutor Note: 易错：(3.29) 的 $\mathcal{K}(a)$ 是待定函数，不是指标；它的作用是把复杂测度吸收进 $\hat\Delta(a)$。]
% [Tutor Note: 物理：(3.31) 的幂次 $2g-2+s$ 正是粘合三穿刺球面所需的个数，由此得到任意亏格曲面的指标。]
% [Tutor Note: 易错：(3.32) 中 $PE$ 的参数分别是 $(-2q/(1-q))$ 与 $(\sqrt q/(1-q))$，指数里 $a$ 的组合不能漏掉 $a^{-2}$ 与 $a^{-1}$。]
% [Tutor Note: 文献：三点函数对角性见 [38]；$T_N$ 指标猜想见 [38][39]；与二维 Yang–Mills 的联系见 [35–37]。]
% [Tutor Note: 物理：$T_{N>2}$ 没有拉氏量，其指标只能靠 TQFT 结构「提升」来猜——这是 class S 指标计算威力的体现。]


\section{大 $N$ 极限下的指标}
% sec4.tex — c08｜原文 p26–30｜§4 大 $N$ 极限下的指标、§4.1 黑洞、Cardy 极限、致谢
% 编号公式 (4.1)–(4.16) 共 16 条；编号环境：习题 4.1；无表格。
% [原文 p.26]
在本节中，我们将计算 $\sN=4$ $\SU(N)$ 超 Yang–Mills 理论在大 $N$ 极限下的超共形指标。其做法是用鞍点近似该积分。同样的鞍点技巧也适用于其他超对称规范理论的大 $N$ 极限 [21,25,40,41]，但这里我们只关注 $\sN=4$ 的情形。当分解为 $\sN=2$ 多重态时，$\sN=4$ 超 Yang–Mills 理论包含一个超多重态和一个矢量多重态，二者都在伴随表示中。单个 $\sN=1$ 矢量多重态的单字母指标为
\[
\begin{aligned}
i(p,q,x,s)&=\frac{-p-q+2pq+(pq)^{\frac13}x^{-1}-(pq)^{\frac23}x}{(1-p)(1-q)}\\
&\quad+\frac{(pq)^{\frac13}x^{\frac12}s-(pq)^{\frac23}x^{-\frac12}s^{-1}+(pq)^{\frac13}x^{\frac12}s^{-1}-(pq)^{\frac23}x^{-\frac12}s}{(1-p)(1-q)}\\
&=\frac{(pq)^{\frac13}(s\sqrt x+\frac{\sqrt x}{s}+\frac1x)-(pq)^{\frac23}(\frac{1}{s\sqrt x}+\frac{s}{\sqrt x}+x)-p-q+2pq}{(1-p)(1-q)} .
\end{aligned}
\tag{4.1}
\]
这里我们引入了逸度 $s$，它对应与 $\sN=4$ 的 $\sN=2$ 子代数对易的荷。它使两个半超多重态带电荷 $\pm1$。要得到 $\SU(N)$ $\sN=4$ 矢量多重态的指标，我们需要把这个指标乘以 $\chi_{\adj}(a_i)$，其中 $a_i$ 是与 $\SU(N)$ 对称性相伴的逸度。有趣的是，单个 $\sN=4$ 矢量的指标满足
\[
1-i(p,q,x,s)=\frac{(1-(pq)^{\frac13}s\sqrt x)(1-(pq)^{\frac13}\frac{\sqrt x}{s})(1-(pq)^{\frac13}\frac1x)}{(1-p)(1-q)} .
\tag{4.2}
\]
这将在后面用到。

规范理论的指标为
\[
\begin{aligned}
\sI&=\frac{1}{N!}\oint\Big(\prod_{j=1}^{N}\frac{\dd a_j}{2\pi\ii}\Big)\Delta(a_i)\,PE\big[i(p,q,x,s)\chi_{\adj}(a_i)\big]\\
&=\frac{1}{N!}\oint\Big(\prod_{j=1}^{N}\frac{\dd a_j}{2\pi\ii}\Big)\exp\Big(-\sum_{n=1}^{\infty}\frac1n\big(1-i(p^n,q^n,x^n,s^n)\big)\sum_{i\neq j}\frac{a_i^n}{a_j^n}\Big) .
\end{aligned}
\tag{4.3}
\]
这类矩阵模型在大 $N$ 极限下已在 [42] 以及 [12] 中被分析。我们遵循他们的分析。引入本征值密度 $\rho(\theta)$ 较为方便，其归一化为 $\int_{-\pi}^{\pi}\dd\theta\,\rho(\theta)=1$。指标于是成为关于 $\rho(\theta)$ 的泛函积分，但要求 $\rho$ 非负。该泛函积分的作用量为
\[
\begin{aligned}
\sI&=\int[\mathcal{D}\rho(\theta)]\exp\Big(N^2\int\dd\theta_1\dd\theta_2\,\rho(\theta_1)V(\theta_1-\theta_2)\rho(\theta_2)\Big)\\
V(\theta)&=-\sum_{n=1}^{\infty}\frac1n\big(1-i(p^n,q^n,x^n,s^n)\cos(n\theta)\big).
\end{aligned}
\tag{4.4}
\]
% [原文 p.27]
这些变量的好处在于，由于作用量中的 $N^2$ 因子，在大 $N$ 时该积分变成鞍点积分。我们进一步用密度模 $\rho_n=\int\dd\theta\,\rho(\theta)e^{in\theta}$ 与势 $V_n=\int\dd\theta\,V(\theta)\cos(n\theta)$ 来简化它。注意 $\rho_{-n}=\rho_n^\dagger$。泛函积分变为
\[
\begin{aligned}
\sI&=\int[\dd\rho_n\,\dd\rho_n^\dagger]\exp\Big(\frac{N^2}{2\pi}\sum_{n=1}^{\infty}\frac1n V_n|\rho_n|^2\Big)\\
V_n&=\frac{2\pi}{n}\big(1-i(p^n,q^n,x^n,s^n)\big).
\end{aligned}
\tag{4.5}
\]
由 (4.2) 可见，对所有 $n\ge1$ 都有 $V_n<0$。这使得所有模 $\rho_{n\ge1}$ 都获得质量。经典位形由 $\rho_0$ 主导。零模的作用量简单地为 $0$，积分由单圈行列式给出，
\[
\begin{aligned}
\sI&=\prod_{n=1}^{\infty}\frac{1}{1-i(p^n,q^n,x^n,s^n)}\\
&=\prod_{n=1}^{\infty}\frac{(1-p^n)(1-q^n)}{(1-(pq)^{\frac n3}s^n x^{\frac n2})(1-(pq)^{\frac n3}s^{-n}x^{\frac n2})(1-(pq)^{\frac n3}x^{-n})} .
\end{aligned}
\tag{4.6}
\]
容易看出，这个闭合形式的指标实际上是某个相对简单的函数的齐次指数。
\[
\begin{aligned}
\sI&=PE[i_G(p,q,x,s)]\\
i_G(p,q,x,s)&=\frac{(pq)^{\frac13}s\sqrt x}{1-(pq)^{\frac13}s\sqrt x}+\frac{(pq)^{\frac13}\frac{\sqrt x}{s}}{1-(pq)^{\frac13}\frac{\sqrt x}{s}}+\frac{(pq)^{\frac13}\frac1x}{1-(pq)^{\frac13}\frac1x}-\frac{p}{1-p}-\frac{q}{1-q} .
\end{aligned}
\tag{4.7}
\]
在大 $N$ 下，$\SU(N)$ $\sN=4$ 超 Yang–Mills 对偶于 $AdS_5\times S^5$ 上的弱耦合 IIB 型弦论。在对偶理论中，只有超引力态受保护（这是因为规范理论变强耦合时，弦模变得非常重）。这意味着 $\sN=4$ SYM 的指标应当计算 $AdS_5\times S^5$ 中超引力场的指标。指标 $\sI$ 的齐次对数，即 $i_G(p,q,x,s)$，确实就是 $AdS_5\times S^5$ 中超引力的单粒子指标 [12]。

\subsection{黑洞呢？}
$AdS_5\times S^5$ 中的超引力也允许超对称黑洞。它们具有 $\mathcal{O}(N^2)$ 的熵。由于我们发现大 $N$ 极限下指标是 $\mathcal{O}(1)$ 的，超对称黑洞的大量态在某种程度上对指标不可见。原因在于，黑洞的玻色态与费米态以这样的方式出现，以致在极大量的态之间发生严重抵消。如果能够避免这些抵消，我们就会找到一个作用量为 $\mathcal{O}(N^2)$ 的鞍点，而不是 $0$。这引出一个自然的问题：指标究竟会不会告诉我们黑洞态的信息？

这个问题的答案取决于玻色态与费米态之间抵消的性质。我们想强调的两种抵消，最好用一个例子
% [原文 p.28]
来说明。考虑指标的单变量极限 $p=q=t^3$、$x=s=1$。在此极限下单字母指标为
\[
i(t)=1-\frac{(1-t^2)^3}{(1-t^3)^2} .
\tag{4.8}
\]
让我们考虑某个大 $N$ 的 $\SU(N)$ 理论指标的 Taylor 级数展开
\[
\sI(t)=\sum_n a_n t^n .
\tag{4.9}
\]
玻色–费米态之间的严重抵消可能极其精细，使得该级数中的每个系数都是 $\mathcal{O}(1)$ 的小量而不是 $\mathcal{O}(N^2)$。在这种情况下，指标将永远无法探测与黑洞熵相伴的大量态简并。但另一种可能是抵消是粗糙的，即级数的系数是 $\mathcal{O}(N^2)$ 的大量，但符号快速振荡。即使在这种情况下，大的简并也难以捕捉，因为指标 $\sI(t)$ 只会看到平均掉的 $\mathcal{O}(1)$ 简并。然而，有一个简单的技巧可以绕开这个问题。若取 $t\to-t$（$t\in\mathbb{R}$），这或许能部分避开玻色–费米抵消，指标本身将是 $\mathcal{O}(N^2)$ 的。若 $a_n$ 的符号振荡周期为 2，则 $t\to-t$ 将完美避开所有抵消。若周期为 3，则我们可能需要取 $t\to t e^{2\pi i/3}$（$t\in\mathbb{R}$）以更好地避开抵消，如此等等。这一讨论的教训是：若玻色–费米态之间的抵消是粗糙的，那么也许可以通过为逸度打开非平庸相位来捕捉 $\mathcal{O}(N^2)$ 的简并增长。

最近，这一可能性已在 [43] 中、在所谓的指标「Cardy 极限」中得到了探索。对该分析而言，用下面描述的更「对称」的方式重新表述 $\sN=4$ 指标更为方便。$\sN=4$ 超共形代数的对易玻色生成元包括三个 R 对称生成元 $r_{i=1,2,3}$、两个转动生成元 $h_{i=1,2}$ 以及一个共形维数 $\Delta$。超荷用其在 Cartan 生成元下的荷 $Q^{r_1,r_2,r_3}_{h_1,h_2}$ 标号。这些荷取值 $\pm\frac12$，使得所有符号的乘积为 $+$。关于 $\mathcal{Q}^{+++}_{--}$ 的超共形指标为
\[
\sI=\Tr\, e^{-\sum_{i=1}^{3}\rho_i r_i}e^{-\sum_{i=1}^{2}\omega_i h_i},\qquad
\rho_1+\rho_2+\rho_3-\omega_1-\omega_2=2\pi i .
\tag{4.10}
\]
\begin{exercise}[4.1]
注意在式 (4.10) 中我们没有包含因子 $(-1)^F$。证明式 (4.10) 仍然定义了一个超共形指标。
提示：所有化学势的周期为 $4\pi i$。
\end{exercise}
% [原文 p.29]
$\sN=4$ 矢量多重态的单字母指标为
\[
i=1-2\,\frac{\sinh\frac{\rho_1}{2}\cdot\sinh\frac{\rho_2}{2}\cdot\sinh\frac{\rho_3}{2}}{\sinh\frac{\omega_1}{2}\cdot\sinh\frac{\omega_2}{2}} .
\tag{4.11}
\]
用这些变量，$\SU(N)$ 规范理论的指标 (4.3) 变为
\[
\sI=\frac{1}{N!}\int_0^{2\pi}\prod_{j=1}^{N}\dd\alpha_j\exp\Big(-\sum_{n=1}^{\infty}\frac2n\frac{\sinh\frac{n\rho_1}{2}\cdot\sinh\frac{n\rho_2}{2}\cdot\sinh\frac{n\rho_3}{2}}{\sinh\frac{n\omega_1}{2}\cdot\sinh\frac{n\omega_2}{2}}\sum_{i\neq j}e^{in\alpha_{ij}}\Big) .
\tag{4.12}
\]
我们不取大 $N$ 极限，而取「Cardy 极限」，即 $\operatorname{Re}(\omega_i)\to0^+$、$\operatorname{Im}(\omega_i)=0$。它被称为 Cardy 极限，是因为它让人联想到取高温极限，只不过在这里这个温度耦合的是角动量而不是能量。注意在此极限下 $\operatorname{Re}(\rho_i)\to0^+$，但由于约束 $\rho_1+\rho_2+\rho_3=\omega_1+\omega_2+2\pi i$ 以及对所有 $i=1,2,3$ 有 $\operatorname{Im}(\rho_i)=\mathcal{O}(1)$，并不需要 $\operatorname{Re}(\rho_i)=0$。在此极限下我们作替换 $\sinh(n\omega_i/2)\to n\omega_i/2$。式 (4.12) 变为
\[
\sI=\frac{1}{N!}\int_0^{2\pi}\prod_{j=1}^{N}\dd\alpha_j\exp\Big(-\frac{1}{\omega_1\omega_2}\sum_{n=1}^{\infty}\frac{2^3}{n^3}\Big(\sinh\frac{n\rho_1}{2}\cdot\sinh\frac{n\rho_2}{2}\cdot\sinh\frac{n\rho_3}{2}\Big)\sum_{i\neq j}e^{in\alpha_{ij}}\Big) .
\tag{4.13}
\]
指数可以进一步用三重对数 $\Li_3(x)\equiv\sum_{n=1}^{\infty}x^n/n^3$ 来化简。
\[
\begin{aligned}
&\sum_{n=1}^{\infty}\frac{2^3}{n^3}\Big(\sinh\frac{n\rho_1}{2}\cdot\sinh\frac{n\rho_2}{2}\cdot\sinh\frac{n\rho_3}{2}\Big)\sum_{i\neq j}e^{in\alpha_{ij}}
=\sum_{s_1,s_2,s_3=\pm1}\sum_{i\neq j}\Li_3\big(e^{\frac{s_i\rho_i}{2}+i\alpha_{ij}}\big)\\
&=\sum_{s_1s_2s_3=1}\sum_{i\neq j}\Big(\Li_3\big(e^{\frac{s_i\rho_i}{2}+i\alpha_{ij}}\big)-\Li_3\big(e^{-\frac{s_i\rho_i}{2}-i\alpha_{ij}}\big)\Big)
\end{aligned}
\tag{4.14}
\]
三重对数具有如下性质
\[
\Li_3(e^x)-\Li_3(e^{-x})=-\frac{x^3}{6}+\frac{\pi i x^2}{2}+\frac{\pi^2 x}{3},\qquad (0<\operatorname{Im}(x)<2\pi,\ \operatorname{Re}(x)\ge0) .
\tag{4.15}
\]
对 $2\pi p<\operatorname{Im}(x)<2\pi(p+1)$，我们需要在右端作平移 $x\to2\pi ip$。经过一番整理，可以注意到作用量泛函的鞍点在 $\alpha_i=\alpha$ 处。该鞍点处作用量的经典值按如下方式标度
\[
\log(\sI)\sim N^2\frac{\rho_1\rho_2\rho_3}{\omega_1\omega_2} .
\tag{4.16}
\]
这正是在 [44] 中被发现、用于产生 $AdS_5$ 中带电 BPS 黑洞正确熵的泛函。同一问题在 [45] 中通过把 Bethe ansatz 方法应用于超共形指标而得到另一种处理。关于有限 $N$ 下指标的数值研究，见 [46,47]。
% [原文 p.30]
\paragraph{致谢}
我们要感谢 DESY 理论物理组，特别是 Elli Pomoni 的盛情款待。我们还要感谢印度人民对基础科学研究的持续支持。

% [Tutor Note: 物理：$\sN=4$ SYM 的指标在大 $N$ 下计数的是 $AdS_5\times S^5$ 超引力态，因为弦模随耦合变强而变得非常重。]
% [Tutor Note: 易错：$s$ 是与 $\sN=2$ 子代数对易的那个 $\U(1)$ 的逸度；两个半超多重态带电荷 $\pm1$，故 (4.1) 中 $s$ 与 $s^{-1}$ 成对出现。]
% [Tutor Note: 易错：(4.1) 的 $x^{\frac12}$ 项对应 $r=1$ 的手征多重态、$x^{-\frac12}$ 项对应 $r=-1$，合并后化成三组因子。]
% [Tutor Note: 物理：(4.2) 的 $1-i$ 恰是 (4.6) 分子的来源；分母的三个因子分别对应 $x^{\frac12}s$、$x^{\frac12}/s$、$x^{-1}$ 的模。]
% [Tutor Note: 易错：(4.3) 中 $\Delta(a_i)$ 与 $PE$ 展开给出 $\sum_{i\neq j}a_i^n/a_j^n$，且不要漏掉前因子 $1/N!$。]
% [Tutor Note: 物理：作用量正比于 $N^2$，大 $N$ 时积分由鞍点主导；$V_n<0$ 使 $\rho_{n\ge1}$ 获得质量，只有零模 $\rho_0$ 无质量。]
% [Tutor Note: 易错：(4.5) 的 $V_n$ 带 $2\pi/n$ 归一化；对 $\rho_n$ 作高斯积分给出 $1/(1-i)$，再取对数即得 $PE$。]
% [Tutor Note: 物理：(4.6) 的闭合形式表明大 $N$ 指标就是 $PE[i_G]$；$i_G$ 与超引力 Kaluza–Klein 谱的单粒子指标一致 [12]。]
% [Tutor Note: 物理：黑洞熵为 $\mathcal{O}(N^2)$ 而指标为 $\mathcal{O}(1)$——典型黑洞态在指标中成对抵消，这正是「指标看不到黑洞」的根源。]
% [Tutor Note: 易错：Cardy 极限不是高温极限；此处温度耦合的是角动量，条件是 $\operatorname{Re}(\omega_i)\to0^+$ 而非 $\beta\to0$。]
% [Tutor Note: 易错：(4.10) 虽无 $(-1)^F$ 仍是指标，因为化学势以 $4\pi i$ 为周期，绕一圈给出的符号可重新解释为 $(-1)^F$（习题 4.1）。]
% [Tutor Note: 易错：(4.12)→(4.13) 只作 $\sinh(n\omega_i/2)\to n\omega_i/2$，并保留 $1/(\omega_1\omega_2)$ 与 $2^3/n^3$。]
% [Tutor Note: 易错：(4.14) 第二行下标 $s_1s_2s_3=1$ 只取乘积为 $+$ 的符号组合，配对 $s\to-s$ 后出现 $\Li_3(x)-\Li_3(1/x)$。]
% [Tutor Note: 文献：(4.16) 与 [44] 的带电 BPS 黑洞熵一致；Bethe ansatz 方法见 [45]；有限 $N$ 数值研究见 [46,47]。]
% [Tutor Note: 易错：原文 (4.4) 第二行括号不平衡，且 $\cos(n\theta)$ 应在因子 $(1-i(\cdots))$ 之外，正确形式见 ERRATA 候选；译文按平衡括号照录。]
% [Tutor Note: 易错：原文 (4.5) 第二行缺右括号，译文补为平衡括号；已登记 ERRATA。]


% ---------------------------- 参考文献 ----------------------------
% refs.tex — c09｜原文 p30–p32｜参考文献 [1]–[47]（照抄不译）
% 本块是全工程唯一允许使用 \section 的块文件（见 BRIEF 与 check_block.py 例外）。
\section{参考文献}
% [原文 p.30]
\begin{enumerate}
\item[{[1]}] S. S. Razamat, \textit{On a modular property of $\sN=2$ superconformal theories in four dimensions}, \textbf{JHEP} 10 (2012) 191 [1208.5056].
\item[{[2]}] A. Gadde, \textit{Modularity of supersymmetric partition functions}, 2004.13490.
\item[{[3]}] V. P. Spiridonov and G. S. Vartanov, \textit{Elliptic hypergeometric integrals and 't Hooft anomaly matching conditions}, \textbf{JHEP} 06 (2012) 016 [1203.5677].
\item[{[4]}] F. Benini, T. Nishioka and M. Yamazaki, \textit{4d Index to 3d Index and 2d TQFT}, \textbf{Phys. Rev.} D86 (2012) 065015 [1109.0283].
\item[{[5]}] S. S. Razamat and B. Willett, \textit{Global Properties of Supersymmetric Theories and the Lens Space}, \textbf{Commun. Math. Phys.} 334 (2015) 661 [1307.4381].
\item[{[6]}] S. S. Razamat and M. Yamazaki, \textit{S-duality and the $N=2$ Lens Space Index}, \textbf{JHEP} 10 (2013) 048 [1306.1543].
\item[{[7]}] A. P. Kels and M. Yamazaki, \textit{Elliptic hypergeometric sum/integral transformations and supersymmetric lens index}, \textbf{SIGMA} 14 (2018) 013 [1704.03159].
\item[{[8]}] L. Rastelli and S. S. Razamat, \textit{The supersymmetric index in four dimensions}, \textbf{J. Phys.} A50 (2017) 443013 [1608.02965].
\item[{[9]}] L. Rastelli and S. S. Razamat, \textit{The Superconformal Index of Theories of Class S}, pp. 261--305. 2016. 1412.7131. 10.1007/978-3-319-18769-3 9.
\item[{[10]}] E. Witten, \textit{Constraints on Supersymmetry Breaking}, \textbf{Nucl. Phys.} B202 (1982) 253.
\item[{[11]}] F. Cooper, A. Khare and U. Sukhatme, \textit{Supersymmetry and quantum mechanics}, \textbf{Phys. Rept.} 251 (1995) 267 [hep-th/9405029].
\item[{[12]}] J. Kinney, J. M. Maldacena, S. Minwalla and S. Raju, \textit{An Index for 4 dimensional super conformal theories}, \textbf{Commun. Math. Phys.} 275 (2007) 209 [hep-th/0510251].
\item[{[13]}] C. Romelsberger, \textit{Counting chiral primaries in $N=1$, $d=4$ superconformal field theories}, \textbf{Nucl. Phys.} B747 (2006) 329 [hep-th/0510060].
\item[{[14]}] C. Romelsberger, \textit{Calculating the Superconformal Index and Seiberg Duality}, 0707.3702.
\item[{[15]}] G. Felder and A. Varchenko, \textit{The elliptic gamma function and $\mathrm{sl}(3,\Z)\ltimes\Z^3$}, \textbf{Advances in Mathematics} 156 (2000) 4476.
\item[{[16]}] K. A. Intriligator and B. Wecht, \textit{The Exact superconformal R symmetry maximizes a}, \textbf{Nucl. Phys.} B 667 (2003) 183 [hep-th/0304128].
\item[{[17]}] D. Kutasov, A. Parnachev and D. A. Sahakyan, \textit{Central charges and $U(1)(R)$ symmetries in $N=1$ superYang-Mills}, \textbf{JHEP} 11 (2003) 013 [hep-th/0308071].
% [原文 p.31]
\item[{[18]}] G. Festuccia and N. Seiberg, \textit{Rigid Supersymmetric Theories in Curved Superspace}, \textbf{JHEP} 06 (2011) 114 [1105.0689].
\item[{[19]}] V. P. Spiridonov, \textit{On the elliptic beta function}, \textbf{Russian Mathematical Surveys} 56 (2001) 185.
\item[{[20]}] N. Seiberg, \textit{Electric - magnetic duality in supersymmetric nonAbelian gauge theories}, \textbf{Nucl. Phys.} B435 (1995) 129 [hep-th/9411149].
\item[{[21]}] F. Dolan and H. Osborn, \textit{Applications of the Superconformal Index for Protected Operators and $q$-Hypergeometric Identities to $N=1$ Dual Theories}, \textbf{Nucl. Phys.} B 818 (2009) 137 [0801.4947].
\item[{[22]}] V. Spiridonov, \textit{Introduction to the theory of elliptic hypergeometric integrals}, 12, 2019, 1912.12971.
\item[{[23]}] F. A. Dolan and H. Osborn, \textit{On short and semi-short representations for four-dimensional superconformal symmetry}, \textbf{Annals Phys.} 307 (2003) 41 [hep-th/0209056].
\item[{[24]}] C. Cordova, T. T. Dumitrescu and K. Intriligator, \textit{Multiplets of Superconformal Symmetry in Diverse Dimensions}, \textbf{JHEP} 03 (2019) 163 [1612.00809].
\item[{[25]}] A. Gadde, L. Rastelli, S. S. Razamat and W. Yan, \textit{On the Superconformal Index of $N=1$ IR Fixed Points: A Holographic Check}, \textbf{JHEP} 03 (2011) 041 [1011.5278].
\item[{[26]}] O. Aharony and M. Evtikhiev, \textit{On four dimensional $N = 3$ superconformal theories}, \textbf{JHEP} 04 (2016) 040 [1512.03524].
\item[{[27]}] I. Garca-Etxebarria and D. Regalado, \textit{$N = 3$ four dimensional field theories}, \textbf{JHEP} 03 (2016) 083 [1512.06434].
\item[{[28]}] O. Aharony and Y. Tachikawa, \textit{S-folds and 4d $N=3$ superconformal field theories}, \textbf{JHEP} 06 (2016) 044 [1602.08638].
\item[{[29]}] T. Bourton, A. Pini and E. Pomoni, \textit{4d $N = 3$ indices via discrete gauging}, \textbf{JHEP} 10 (2018) 131 [1804.05396].
\item[{[30]}] A. Gadde, L. Rastelli, S. S. Razamat and W. Yan, \textit{Gauge Theories and Macdonald Polynomials}, \textbf{Commun. Math. Phys.} 319 (2013) 147 [1110.3740].
\item[{[31]}] D. Gaiotto, \textit{$N=2$ dualities}, \textbf{JHEP} 08 (2012) 034 [0904.2715].
\item[{[32]}] N. Seiberg and E. Witten, \textit{Monopoles, duality and chiral symmetry breaking in $N=2$ supersymmetric QCD}, \textbf{Nucl. Phys.} B431 (1994) 484 [hep-th/9408099].
\item[{[33]}] F. J. van de Bult, \textit{An elliptic hypergeometric integral with $w(f_4)$ symmetry}, \textbf{The Ramanujan Journal} 25 (2011) 1.
\item[{[34]}] A. Gadde, E. Pomoni, L. Rastelli and S. S. Razamat, \textit{S-duality and 2d Topological QFT}, \textbf{JHEP} 03 (2010) 032 [0910.2225].
\item[{[35]}] E. Witten, \textit{On quantum gauge theories in two-dimensions}, \textbf{Commun. Math. Phys.} 141 (1991) 153.
\item[{[36]}] E. Witten, \textit{Two-dimensional gauge theories revisited}, \textbf{J. Geom. Phys.} 9 (1992) 303 [hep-th/9204083].
\item[{[37]}] S. Cordes, G. W. Moore and S. Ramgoolam, \textit{Lectures on 2-d Yang-Mills theory, equivariant cohomology and topological field theories}, \textbf{Nucl. Phys. Proc. Suppl.} 41 (1995) 184 [hep-th/9411210].
% [原文 p.32]
\item[{[38]}] A. Gadde, L. Rastelli, S. S. Razamat and W. Yan, \textit{The 4d Superconformal Index from $q$-deformed 2d Yang-Mills}, \textbf{Phys. Rev. Lett.} 106 (2011) 241602 [1104.3850].
\item[{[39]}] D. Gaiotto, L. Rastelli and S. S. Razamat, \textit{Bootstrapping the superconformal index with surface defects}, \textbf{JHEP} 01 (2013) 022 [1207.3577].
\item[{[40]}] Y. Nakayama, \textit{Index for orbifold quiver gauge theories}, \textbf{Phys. Lett.} B 636 (2006) 132 [hep-th/0512280].
\item[{[41]}] Y. Nakayama, \textit{Index for supergravity on AdS(5) x T**1,1 and conifold gauge theory}, \textbf{Nucl. Phys.} B 755 (2006) 295 [hep-th/0602284].
\item[{[42]}] O. Aharony, J. Marsano, S. Minwalla, K. Papadodimas and M. Van Raamsdonk, \textit{The Hagedorn - deconfinement phase transition in weakly coupled large $N$ gauge theories}, \textbf{Adv. Theor. Math. Phys.} 8 (2004) 603 [hep-th/0310285].
\item[{[43]}] S. Choi, J. Kim, S. Kim and J. Nahmgoong, \textit{Large AdS black holes from QFT}, 1810.12067.
\item[{[44]}] S. M. Hosseini, K. Hristov and A. Zaffaroni, \textit{An extremization principle for the entropy of rotating BPS black holes in $AdS_5$}, \textbf{JHEP} 07 (2017) 106 [1705.05383].
\item[{[45]}] F. Benini and P. Milan, \textit{Black holes in 4d $\sN = 4$ Super-Yang-Mills}, 1812.09613.
\item[{[46]}] P. Agarwal, S. Choi, J. Kim, S. Kim and J. Nahmgoong, \textit{AdS black holes and finite $N$ indices}, 2005.11240.
\item[{[47]}] S. Murthy, \textit{The growth of the $\frac{1}{16}$-BPS index in 4d $\sN = 4$ SYM}, 2005.10843.
\end{enumerate}

% [Tutor Note: 本块按 BRIEF 要求「照抄不译」，仅把原文的斜体标题、粗体刊名与数学记号复原；文献内容一字未改。]
% [Tutor Note: [27] 作者名原文即印作 “Garca-Etxebarria”（缺 í），照录，未擅改。]
% [Tutor Note: [15] 页码原文印作 “4476”，疑为 44–76 的排版粘连，照录并见 ERRATA 候选。]
% [Tutor Note: [41] 原文把 $T^{1,1}$ 印成 “T**1,1”，照录（见 ERRATA 候选）。]


\end{document}
